
\documentclass{article}
\usepackage{colortbl}
\usepackage{makecell}
\usepackage{multirow}
\usepackage{supertabular}

\begin{document}

\newcounter{utterance}

\centering \large Interaction Transcript for game `clean\_up', experiment `2\_hard\_5obj\_en', episode 2 with Llama{-}3.3{-}70B{-}Instruct{-}t0.0.
\vspace{24pt}

{ \footnotesize  \setcounter{utterance}{1}
\setlength{\tabcolsep}{0pt}
\begin{supertabular}{c@{$\;$}|p{.15\linewidth}@{}p{.15\linewidth}p{.15\linewidth}p{.15\linewidth}p{.15\linewidth}p{.15\linewidth}}
    \# & $\;$A & \multicolumn{4}{c}{Game Master} & $\;\:$B\\
    \hline

    \theutterance \stepcounter{utterance}  
    & & \multicolumn{4}{p{0.6\linewidth}}{
        \cellcolor[rgb]{0.9,0.9,0.9}{
            \makecell[{{p{\linewidth}}}]{
                \texttt{\tiny{[P1$\langle$GM]}}
                \texttt{I am your game master, and you are playing a collaborative game with the following grid as game board:} \\
\\ 
\texttt{```} \\
\texttt{ 1234567} \\
\texttt{╔═╤═╤═╤═╗} \\
\texttt{║T│◌│◌│◌║ 1} \\
\texttt{╟─┼─┤I│◌║ 2} \\
\texttt{║W│◌│◌│H║ 3} \\
\texttt{║◌│◌├─┴─╢ 4} \\
\texttt{╟─┴─┤◌◌◌║ 5} \\
\texttt{║◌C◌└─┬─╢ 6} \\
\texttt{║◌◌◌◌◌│◌║ 7} \\
\texttt{╚═════╧═╝} \\
\\ 
\texttt{```} \\
\\ 
\texttt{* The upper edge displays x{-}coordinates increasing to the right, and the right edge y{-}coordinates increasing downward.} \\
\texttt{* The following objects are randomly placed on your grid: 'W', 'I', 'T', 'C', 'H'.} \\
\\ 
\\ 
\texttt{The other player sees a variation of the game board, where the objects are placed at different random locations. You cannot see the other player's board, and they cannot see yours.} \\
\\ 
\\ 
\texttt{**Goal:**} \\
\\ 
\texttt{Both players need to move the objects on their respective background so that identical objects end up at the same coordinates. You have to communicate with the other player to agree upon a common goal configuration.} \\
\\ 
\\ 
\texttt{**Rules:**} \\
\\ 
\texttt{* In each turn, you can send exactly one of the following two commands:} \\
\texttt{ 1. `SAY: <MESSAGE>`: to send a message (everything up until the next line break) to the other player. I will forward it to your partner.} \\
\texttt{2. `MOVE: <OBJECT>, (<X>, <Y>)`: to move an object to a new position, where `<X>` is the column and `<Y>` is the row. I will inform you if your move was successful or not.} \\
\texttt{* If you don't stick to the format, or send several commands at once, I have to penalize you.} \\
\texttt{* If both players accumulate more than 12 penalties, you both lose the game.} \\
\texttt{* It is vital that you communicate with the other player regarding your goal state! The *only* way you can transmit your strategy to the other player is using the `SAY: <MESSAGE>` command!} \\
\\ 
\\ 
\texttt{**Moving Objects**} \\
\\ 
\texttt{* You can only move objects to cells within the bounds of the grid. The target cell must be empty, i.e., it must only contain the symbol '◌'.} \\
\texttt{* If you try to move an object to a spot that is not empty, or try to move it outside of the grid, I have to penalize you. You get another try.} \\
\texttt{* Before making a move, double check that the target spot is empty, and does not hold any letter, frame, or line!} \\
\\ 
\\ 
\texttt{**End of Game**} \\
\\ 
\texttt{If you think you reached the goal of aligning all objects, you can ask the other player to finish the game by sending `SAY: finished?`. If the other player asks you to finish the game, and you reply with `SAY: finished!`, the game will end.} \\
\\ 
\texttt{Both players win if the game ends within 20 rounds, where one round is defined as two players each sending a valid command.} \\
\\ 
\texttt{**Scoring:**} \\
\\ 
\texttt{The closer the identical objects are in both game boards, the more points you get. Penalties reduce your points. Can you beat the record?} \\
\\ 
\\ 
\texttt{Please send a message to the other player to start the game!} \\
            }
        }
    }
    & \\ \\

    \theutterance \stepcounter{utterance}  
    & \multicolumn{4}{p{0.6\linewidth}}{
        \cellcolor[rgb]{0.9,0.9,0.9}{
            \makecell[{{p{\linewidth}}}]{
                \texttt{\tiny{[P1$\rangle$GM]}}
                \texttt{SAY: Hello, let's start by describing our initial boards. I'll describe mine, and you describe yours. My board has the following objects at these coordinates: T(1,1), I(4,2), W(1,3), H(4,3), C(2,6). Please respond with the coordinates of the objects on your board.} \\
            }
        }
    }
    & & \\ \\

    \theutterance \stepcounter{utterance}  
    & & & \multicolumn{2}{p{0.3\linewidth}}{
        \cellcolor[rgb]{0.9,0.9,0.9}{
            \makecell[{{p{\linewidth}}}]{
                \texttt{\tiny{[GM$|$GM]}}
                \texttt{SAY: Hello, let's start by describing our initial boards. I'll describe mine, and you describe yours. My board has the following objects at these coordinates: T(1,1), I(4,2), W(1,3), H(4,3), C(2,6). Please respond with the coordinates of the objects on your board.} \\
            }
        }
    }
    & & \\ \\

    \theutterance \stepcounter{utterance}  
    & & \multicolumn{4}{p{0.6\linewidth}}{
        \cellcolor[rgb]{0.9,0.9,0.9}{
            \makecell[{{p{\linewidth}}}]{
                \texttt{\tiny{[GM$\rangle$P2]}}
                \texttt{I am your game master, and you are playing a collaborative game with the following grid as game board:} \\
\\ 
\texttt{```} \\
\texttt{ 1234567} \\
\texttt{╔═╤═╤═╤═╗} \\
\texttt{║◌│◌│◌│◌║ 1} \\
\texttt{╟─┼─┤◌│◌║ 2} \\
\texttt{║◌│W│◌│◌║ 3} \\
\texttt{║◌│C├─┴─╢ 4} \\
\texttt{╟─┴─┤I◌◌║ 5} \\
\texttt{║◌H◌└─┬─╢ 6} \\
\texttt{║◌◌T◌◌│◌║ 7} \\
\texttt{╚═════╧═╝} \\
\\ 
\texttt{```} \\
\\ 
\texttt{* The upper edge displays x{-}coordinates increasing to the right, and the right edge y{-}coordinates increasing downward.} \\
\texttt{* The following objects are randomly placed on your grid: 'W', 'I', 'T', 'C', 'H'.} \\
\\ 
\\ 
\texttt{The other player sees a variation of the game board, where the objects are placed at different random locations. You cannot see the other player's board, and they cannot see yours.} \\
\\ 
\\ 
\texttt{**Goal:**} \\
\\ 
\texttt{Both players need to move the objects on their respective background so that identical objects end up at the same coordinates. You have to communicate with the other player to agree upon a common goal configuration.} \\
\\ 
\\ 
\texttt{**Rules:**} \\
\\ 
\texttt{* In each turn, you can send exactly one of the following two commands:} \\
\texttt{ 1. `SAY: <MESSAGE>`: to send a message (everything up until the next line break) to the other player. I will forward it to your partner.} \\
\texttt{2. `MOVE: <OBJECT>, (<X>, <Y>)`: to move an object to a new position, where `<X>` is the column and `<Y>` is the row. I will inform you if your move was successful or not.} \\
\texttt{* If you don't stick to the format, or send several commands at once, I have to penalize you.} \\
\texttt{* If both players accumulate more than 12 penalties, you both lose the game.} \\
\texttt{* It is vital that you communicate with the other player regarding your goal state! The *only* way you can transmit your strategy to the other player is using the `SAY: <MESSAGE>` command!} \\
\\ 
\\ 
\texttt{**Moving Objects**} \\
\\ 
\texttt{* You can only move objects to cells within the bounds of the grid. The target cell must be empty, i.e., it must only contain the symbol '◌'.} \\
\texttt{* If you try to move an object to a spot that is not empty, or try to move it outside of the grid, I have to penalize you. You get another try.} \\
\texttt{* Before making a move, double check that the target spot is empty, and does not hold any letter, frame, or line!} \\
\\ 
\\ 
\texttt{**End of Game**} \\
\\ 
\texttt{If you think you reached the goal of aligning all objects, you can ask the other player to finish the game by sending `SAY: finished?`. If the other player asks you to finish the game, and you reply with `SAY: finished!`, the game will end.} \\
\\ 
\texttt{Both players win if the game ends within 20 rounds, where one round is defined as two players each sending a valid command.} \\
\\ 
\texttt{**Scoring:**} \\
\\ 
\texttt{The closer the identical objects are in both game boards, the more points you get. Penalties reduce your points. Can you beat the record?} \\
\\ 
\texttt{The other player started the game by sending this message:} \\
\\ 
\texttt{"Hello, let's start by describing our initial boards. I'll describe mine, and you describe yours. My board has the following objects at these coordinates: T(1,1), I(4,2), W(1,3), H(4,3), C(2,6). Please respond with the coordinates of the objects on your board."} \\
\\ 
\texttt{What is your first command?} \\
            }
        }
    }
    & \\ \\

    \theutterance \stepcounter{utterance}  
    & & & \multicolumn{4}{p{0.6\linewidth}}{
        \cellcolor[rgb]{0.9,0.9,0.9}{
            \makecell[{{p{\linewidth}}}]{
                \texttt{\tiny{[GM$\langle$P2]}}
                \texttt{My first command is to respond to the other player's message with the coordinates of the objects on my board.} \\
\\ 
\texttt{SAY: Hello! My board has the following objects at these coordinates: W(2,3), I(1,5), T(3,7), C(2,4), H(1,6). Let's discuss how to align them.} \\
            }
        }
    }
     \\ \\

    \theutterance \stepcounter{utterance}  
    & & & \multicolumn{2}{p{0.3\linewidth}}{
        \cellcolor[rgb]{0.9,0.9,0.9}{
            \makecell[{{p{\linewidth}}}]{
                \texttt{\tiny{[GM$|$GM]}}
                \texttt{My first command is to respond to the other player's message with the coordinates of the objects on my board.} \\
\\ 
\texttt{SAY: Hello! My board has the following objects at these coordinates: W(2,3), I(1,5), T(3,7), C(2,4), H(1,6). Let's discuss how to align them.} \\
            }
        }
    }
    & & \\ \\

    \theutterance \stepcounter{utterance}  
    & & & \multicolumn{2}{p{0.3\linewidth}}{
        \cellcolor[rgb]{0.9,0.9,0.9}{
            \makecell[{{p{\linewidth}}}]{
                \texttt{\tiny{[GM$|$GM]}}
                \texttt{Invalid format: head is not empty: 'My first command is to respond to the other player's message with the coordinates of the objects on my board.} \\
\\ 
\texttt{'} \\
            }
        }
    }
    & & \\ \\

    \theutterance \stepcounter{utterance}  
    & & \multicolumn{4}{p{0.6\linewidth}}{
        \cellcolor[rgb]{0.9,0.9,0.9}{
            \makecell[{{p{\linewidth}}}]{
                \texttt{\tiny{[GM$\rangle$P2]}}
                \texttt{Penalty: Your message must not contain anything before the command!} \\
\\ 
\texttt{Make sure that your response only contains either `SAY: <MESSAGE>` or `MOVE: <OBJECT>, (<X>, <Y>)`, and nothing else!} \\
\\ 
\texttt{You have collectively accumulated 1 of 12 penalties.} \\
\texttt{Please try again!} \\
            }
        }
    }
    & \\ \\

    \theutterance \stepcounter{utterance}  
    & & & \multicolumn{4}{p{0.6\linewidth}}{
        \cellcolor[rgb]{0.9,0.9,0.9}{
            \makecell[{{p{\linewidth}}}]{
                \texttt{\tiny{[GM$\langle$P2]}}
                \texttt{SAY: My board has the following objects at these coordinates: W(2,3), I(1,5), T(3,7), C(2,4), H(1,6). Let's discuss how to align them.} \\
            }
        }
    }
     \\ \\

    \theutterance \stepcounter{utterance}  
    & & & \multicolumn{2}{p{0.3\linewidth}}{
        \cellcolor[rgb]{0.9,0.9,0.9}{
            \makecell[{{p{\linewidth}}}]{
                \texttt{\tiny{[GM$|$GM]}}
                \texttt{SAY: My board has the following objects at these coordinates: W(2,3), I(1,5), T(3,7), C(2,4), H(1,6). Let's discuss how to align them.} \\
            }
        }
    }
    & & \\ \\

    \theutterance \stepcounter{utterance}  
    & & \multicolumn{4}{p{0.6\linewidth}}{
        \cellcolor[rgb]{0.9,0.9,0.9}{
            \makecell[{{p{\linewidth}}}]{
                \texttt{\tiny{[P1$\langle$GM]}}
                \texttt{Your message has been relayed to the other player.} \\
\\ 
\\ 
\texttt{You are currently playing round 1 of maximum 20.} \\
\texttt{You have collectively accumulated 1 of 12 penalties.} \\
\texttt{The other player sent this message:} \\
\\ 
\\ 
\texttt{"My board has the following objects at these coordinates: W(2,3), I(1,5), T(3,7), C(2,4), H(1,6). Let's discuss how to align them."} \\
\\ 
\\ 
\texttt{What is your next command?} \\
            }
        }
    }
    & \\ \\

    \theutterance \stepcounter{utterance}  
    & \multicolumn{4}{p{0.6\linewidth}}{
        \cellcolor[rgb]{0.9,0.9,0.9}{
            \makecell[{{p{\linewidth}}}]{
                \texttt{\tiny{[P1$\rangle$GM]}}
                \texttt{SAY: It seems our boards have different initial placements. To align the objects, let's start by agreeing on a common position for one object. How about we try to move all the 'T's to the same coordinates? I propose we aim for T(3,7) since that's where your 'T' is already placed. What do you think? Should we prioritize moving our 'T' to (3,7) or do you have a different suggestion?} \\
            }
        }
    }
    & & \\ \\

    \theutterance \stepcounter{utterance}  
    & & & \multicolumn{2}{p{0.3\linewidth}}{
        \cellcolor[rgb]{0.9,0.9,0.9}{
            \makecell[{{p{\linewidth}}}]{
                \texttt{\tiny{[GM$|$GM]}}
                \texttt{SAY: It seems our boards have different initial placements. To align the objects, let's start by agreeing on a common position for one object. How about we try to move all the 'T's to the same coordinates? I propose we aim for T(3,7) since that's where your 'T' is already placed. What do you think? Should we prioritize moving our 'T' to (3,7) or do you have a different suggestion?} \\
            }
        }
    }
    & & \\ \\

    \theutterance \stepcounter{utterance}  
    & & \multicolumn{4}{p{0.6\linewidth}}{
        \cellcolor[rgb]{0.9,0.9,0.9}{
            \makecell[{{p{\linewidth}}}]{
                \texttt{\tiny{[GM$\rangle$P2]}}
                \texttt{Your message has been relayed to the other player.} \\
\\ 
\\ 
\texttt{You are currently playing round 2 of maximum 20.} \\
\texttt{You have collectively accumulated 1 of 12 penalties.} \\
\texttt{The other player sent this message:} \\
\\ 
\\ 
\texttt{"It seems our boards have different initial placements. To align the objects, let's start by agreeing on a common position for one object. How about we try to move all the 'T's to the same coordinates? I propose we aim for T(3,7) since that's where your 'T' is already placed. What do you think? Should we prioritize moving our 'T' to (3,7) or do you have a different suggestion?"} \\
\\ 
\\ 
\texttt{What is your next command?} \\
            }
        }
    }
    & \\ \\

    \theutterance \stepcounter{utterance}  
    & & & \multicolumn{4}{p{0.6\linewidth}}{
        \cellcolor[rgb]{0.9,0.9,0.9}{
            \makecell[{{p{\linewidth}}}]{
                \texttt{\tiny{[GM$\langle$P2]}}
                \texttt{SAY: That sounds like a good plan. I agree to aim for T(3,7). Since my 'T' is already at (3,7), I won't need to move it. What's your next step to move your 'T' to (3,7)? Should we discuss the positions of the other objects after this?} \\
            }
        }
    }
     \\ \\

    \theutterance \stepcounter{utterance}  
    & & & \multicolumn{2}{p{0.3\linewidth}}{
        \cellcolor[rgb]{0.9,0.9,0.9}{
            \makecell[{{p{\linewidth}}}]{
                \texttt{\tiny{[GM$|$GM]}}
                \texttt{SAY: That sounds like a good plan. I agree to aim for T(3,7). Since my 'T' is already at (3,7), I won't need to move it. What's your next step to move your 'T' to (3,7)? Should we discuss the positions of the other objects after this?} \\
            }
        }
    }
    & & \\ \\

    \theutterance \stepcounter{utterance}  
    & & \multicolumn{4}{p{0.6\linewidth}}{
        \cellcolor[rgb]{0.9,0.9,0.9}{
            \makecell[{{p{\linewidth}}}]{
                \texttt{\tiny{[P1$\langle$GM]}}
                \texttt{Your message has been relayed to the other player.} \\
\\ 
\\ 
\texttt{You are currently playing round 2 of maximum 20.} \\
\texttt{You have collectively accumulated 1 of 12 penalties.} \\
\texttt{The other player sent this message:} \\
\\ 
\\ 
\texttt{"That sounds like a good plan. I agree to aim for T(3,7). Since my 'T' is already at (3,7), I won't need to move it. What's your next step to move your 'T' to (3,7)? Should we discuss the positions of the other objects after this?"} \\
\\ 
\\ 
\texttt{What is your next command?} \\
            }
        }
    }
    & \\ \\

    \theutterance \stepcounter{utterance}  
    & \multicolumn{4}{p{0.6\linewidth}}{
        \cellcolor[rgb]{0.9,0.9,0.9}{
            \makecell[{{p{\linewidth}}}]{
                \texttt{\tiny{[P1$\rangle$GM]}}
                \texttt{SAY: Great, let's focus on moving my 'T' to (3,7) first. Before I make the move, I want to confirm that the path is clear. My 'T' is currently at (1,1), so I'll need to move it 2 columns to the right and 6 rows down. After my 'T' is in place, we can discuss the positions of the other objects. How about we aim to place the 'I's, 'W's, 'C's, and 'H's in a specific order, like alphabetical order, to simplify the process? What are your thoughts on this approach?} \\
            }
        }
    }
    & & \\ \\

    \theutterance \stepcounter{utterance}  
    & & & \multicolumn{2}{p{0.3\linewidth}}{
        \cellcolor[rgb]{0.9,0.9,0.9}{
            \makecell[{{p{\linewidth}}}]{
                \texttt{\tiny{[GM$|$GM]}}
                \texttt{SAY: Great, let's focus on moving my 'T' to (3,7) first. Before I make the move, I want to confirm that the path is clear. My 'T' is currently at (1,1), so I'll need to move it 2 columns to the right and 6 rows down. After my 'T' is in place, we can discuss the positions of the other objects. How about we aim to place the 'I's, 'W's, 'C's, and 'H's in a specific order, like alphabetical order, to simplify the process? What are your thoughts on this approach?} \\
            }
        }
    }
    & & \\ \\

    \theutterance \stepcounter{utterance}  
    & & \multicolumn{4}{p{0.6\linewidth}}{
        \cellcolor[rgb]{0.9,0.9,0.9}{
            \makecell[{{p{\linewidth}}}]{
                \texttt{\tiny{[GM$\rangle$P2]}}
                \texttt{Your message has been relayed to the other player.} \\
\\ 
\\ 
\texttt{You are currently playing round 3 of maximum 20.} \\
\texttt{You have collectively accumulated 1 of 12 penalties.} \\
\texttt{The other player sent this message:} \\
\\ 
\\ 
\texttt{"Great, let's focus on moving my 'T' to (3,7) first. Before I make the move, I want to confirm that the path is clear. My 'T' is currently at (1,1), so I'll need to move it 2 columns to the right and 6 rows down. After my 'T' is in place, we can discuss the positions of the other objects. How about we aim to place the 'I's, 'W's, 'C's, and 'H's in a specific order, like alphabetical order, to simplify the process? What are your thoughts on this approach?"} \\
\\ 
\\ 
\texttt{What is your next command?} \\
            }
        }
    }
    & \\ \\

    \theutterance \stepcounter{utterance}  
    & & & \multicolumn{4}{p{0.6\linewidth}}{
        \cellcolor[rgb]{0.9,0.9,0.9}{
            \makecell[{{p{\linewidth}}}]{
                \texttt{\tiny{[GM$\langle$P2]}}
                \texttt{SAY: Alphabetical order sounds like a good approach to simplify the process. It would be easier to keep track of the objects' positions. After your 'T' is in place, we can discuss the specific coordinates for the other objects in alphabetical order. I'm waiting for your 'T' to be moved to (3,7) before we proceed.} \\
            }
        }
    }
     \\ \\

    \theutterance \stepcounter{utterance}  
    & & & \multicolumn{2}{p{0.3\linewidth}}{
        \cellcolor[rgb]{0.9,0.9,0.9}{
            \makecell[{{p{\linewidth}}}]{
                \texttt{\tiny{[GM$|$GM]}}
                \texttt{SAY: Alphabetical order sounds like a good approach to simplify the process. It would be easier to keep track of the objects' positions. After your 'T' is in place, we can discuss the specific coordinates for the other objects in alphabetical order. I'm waiting for your 'T' to be moved to (3,7) before we proceed.} \\
            }
        }
    }
    & & \\ \\

    \theutterance \stepcounter{utterance}  
    & & \multicolumn{4}{p{0.6\linewidth}}{
        \cellcolor[rgb]{0.9,0.9,0.9}{
            \makecell[{{p{\linewidth}}}]{
                \texttt{\tiny{[P1$\langle$GM]}}
                \texttt{Your message has been relayed to the other player.} \\
\\ 
\\ 
\texttt{You are currently playing round 3 of maximum 20.} \\
\texttt{You have collectively accumulated 1 of 12 penalties.} \\
\texttt{The other player sent this message:} \\
\\ 
\\ 
\texttt{"Alphabetical order sounds like a good approach to simplify the process. It would be easier to keep track of the objects' positions. After your 'T' is in place, we can discuss the specific coordinates for the other objects in alphabetical order. I'm waiting for your 'T' to be moved to (3,7) before we proceed."} \\
\\ 
\\ 
\texttt{What is your next command?} \\
            }
        }
    }
    & \\ \\

    \theutterance \stepcounter{utterance}  
    & \multicolumn{4}{p{0.6\linewidth}}{
        \cellcolor[rgb]{0.9,0.9,0.9}{
            \makecell[{{p{\linewidth}}}]{
                \texttt{\tiny{[P1$\rangle$GM]}}
                \texttt{MOVE: T, (3, 7)} \\
            }
        }
    }
    & & \\ \\

    \theutterance \stepcounter{utterance}  
    & & & \multicolumn{2}{p{0.3\linewidth}}{
        \cellcolor[rgb]{0.9,0.9,0.9}{
            \makecell[{{p{\linewidth}}}]{
                \texttt{\tiny{[GM$|$GM]}}
                \texttt{MOVE: T, (3, 7)} \\
            }
        }
    }
    & & \\ \\

    \theutterance \stepcounter{utterance}  
    & & & \multicolumn{2}{p{0.3\linewidth}}{
        \cellcolor[rgb]{0.9,0.9,0.9}{
            \makecell[{{p{\linewidth}}}]{
                \texttt{\tiny{[GM$|$GM]}}
                \texttt{T moved to (3, 7)} \\
            }
        }
    }
    & & \\ \\

    \theutterance \stepcounter{utterance}  
    & & \multicolumn{4}{p{0.6\linewidth}}{
        \cellcolor[rgb]{0.9,0.9,0.9}{
            \makecell[{{p{\linewidth}}}]{
                \texttt{\tiny{[GM$\rangle$P2]}}
                \texttt{Your message has been relayed to the other player.} \\
\\ 
\\ 
\texttt{You are currently playing round 4 of maximum 20.} \\
\texttt{You have collectively accumulated 1 of 12 penalties.} \\
\texttt{The other player moved an object on their grid.} \\
\\ 
\\ 
\texttt{What is your next command?} \\
            }
        }
    }
    & \\ \\

    \theutterance \stepcounter{utterance}  
    & & & \multicolumn{4}{p{0.6\linewidth}}{
        \cellcolor[rgb]{0.9,0.9,0.9}{
            \makecell[{{p{\linewidth}}}]{
                \texttt{\tiny{[GM$\langle$P2]}}
                \texttt{SAY: Now that you've moved your object, can you confirm if your 'T' is indeed at (3,7)? If so, we can start discussing the positions of the other objects in alphabetical order, starting with 'C'. What are your thoughts on where we should place the 'C's?} \\
            }
        }
    }
     \\ \\

    \theutterance \stepcounter{utterance}  
    & & & \multicolumn{2}{p{0.3\linewidth}}{
        \cellcolor[rgb]{0.9,0.9,0.9}{
            \makecell[{{p{\linewidth}}}]{
                \texttt{\tiny{[GM$|$GM]}}
                \texttt{SAY: Now that you've moved your object, can you confirm if your 'T' is indeed at (3,7)? If so, we can start discussing the positions of the other objects in alphabetical order, starting with 'C'. What are your thoughts on where we should place the 'C's?} \\
            }
        }
    }
    & & \\ \\

    \theutterance \stepcounter{utterance}  
    & & \multicolumn{4}{p{0.6\linewidth}}{
        \cellcolor[rgb]{0.9,0.9,0.9}{
            \makecell[{{p{\linewidth}}}]{
                \texttt{\tiny{[P1$\langle$GM]}}
                \texttt{Moved 'T' to (3,7) successfully. Your updated grid looks like this:} \\
\texttt{```} \\
\texttt{ 1234567} \\
\texttt{╔═╤═╤═╤═╗} \\
\texttt{║◌│◌│◌│◌║ 1} \\
\texttt{╟─┼─┤I│◌║ 2} \\
\texttt{║W│◌│◌│H║ 3} \\
\texttt{║◌│◌├─┴─╢ 4} \\
\texttt{╟─┴─┤◌◌◌║ 5} \\
\texttt{║◌C◌└─┬─╢ 6} \\
\texttt{║◌◌T◌◌│◌║ 7} \\
\texttt{╚═════╧═╝} \\
\texttt{```} \\
\texttt{You are currently playing round 4 of maximum 20.} \\
\texttt{You have collectively accumulated 1 of 12 penalties.} \\
\texttt{The other player sent this message:} \\
\\ 
\\ 
\texttt{"Now that you've moved your object, can you confirm if your 'T' is indeed at (3,7)? If so, we can start discussing the positions of the other objects in alphabetical order, starting with 'C'. What are your thoughts on where we should place the 'C's?"} \\
\\ 
\\ 
\texttt{What is your next command?} \\
            }
        }
    }
    & \\ \\

    \theutterance \stepcounter{utterance}  
    & \multicolumn{4}{p{0.6\linewidth}}{
        \cellcolor[rgb]{0.9,0.9,0.9}{
            \makecell[{{p{\linewidth}}}]{
                \texttt{\tiny{[P1$\rangle$GM]}}
                \texttt{SAY: Yes, I confirm that my 'T' is now at (3,7). Let's discuss the 'C's. Since we're aiming for a consistent placement, how about we place the 'C's at the same row, but for now, let's focus on finding a suitable column? Perhaps we could place the 'C's at a column that is relatively close to the center or has the least number of objects currently? My 'C' is currently at (2,6). What's the position of your 'C', and do you have any suggestions for the 'C's coordinates?} \\
            }
        }
    }
    & & \\ \\

    \theutterance \stepcounter{utterance}  
    & & & \multicolumn{2}{p{0.3\linewidth}}{
        \cellcolor[rgb]{0.9,0.9,0.9}{
            \makecell[{{p{\linewidth}}}]{
                \texttt{\tiny{[GM$|$GM]}}
                \texttt{SAY: Yes, I confirm that my 'T' is now at (3,7). Let's discuss the 'C's. Since we're aiming for a consistent placement, how about we place the 'C's at the same row, but for now, let's focus on finding a suitable column? Perhaps we could place the 'C's at a column that is relatively close to the center or has the least number of objects currently? My 'C' is currently at (2,6). What's the position of your 'C', and do you have any suggestions for the 'C's coordinates?} \\
            }
        }
    }
    & & \\ \\

    \theutterance \stepcounter{utterance}  
    & & \multicolumn{4}{p{0.6\linewidth}}{
        \cellcolor[rgb]{0.9,0.9,0.9}{
            \makecell[{{p{\linewidth}}}]{
                \texttt{\tiny{[GM$\rangle$P2]}}
                \texttt{Your message has been relayed to the other player.} \\
\\ 
\\ 
\texttt{You are currently playing round 5 of maximum 20.} \\
\texttt{You have collectively accumulated 1 of 12 penalties.} \\
\texttt{The other player sent this message:} \\
\\ 
\\ 
\texttt{"Yes, I confirm that my 'T' is now at (3,7). Let's discuss the 'C's. Since we're aiming for a consistent placement, how about we place the 'C's at the same row, but for now, let's focus on finding a suitable column? Perhaps we could place the 'C's at a column that is relatively close to the center or has the least number of objects currently? My 'C' is currently at (2,6). What's the position of your 'C', and do you have any suggestions for the 'C's coordinates?"} \\
\\ 
\\ 
\texttt{What is your next command?} \\
            }
        }
    }
    & \\ \\

    \theutterance \stepcounter{utterance}  
    & & & \multicolumn{4}{p{0.6\linewidth}}{
        \cellcolor[rgb]{0.9,0.9,0.9}{
            \makecell[{{p{\linewidth}}}]{
                \texttt{\tiny{[GM$\langle$P2]}}
                \texttt{SAY: My 'C' is currently at (2,4). Considering the grid and the current positions, I think placing the 'C's at column 2 could be a good idea, since both our 'C's are already at that column. We could aim to place them at the same row, perhaps at a row that is relatively empty or close to the center. What do you think about placing the 'C's at row 4 or 5?} \\
            }
        }
    }
     \\ \\

    \theutterance \stepcounter{utterance}  
    & & & \multicolumn{2}{p{0.3\linewidth}}{
        \cellcolor[rgb]{0.9,0.9,0.9}{
            \makecell[{{p{\linewidth}}}]{
                \texttt{\tiny{[GM$|$GM]}}
                \texttt{SAY: My 'C' is currently at (2,4). Considering the grid and the current positions, I think placing the 'C's at column 2 could be a good idea, since both our 'C's are already at that column. We could aim to place them at the same row, perhaps at a row that is relatively empty or close to the center. What do you think about placing the 'C's at row 4 or 5?} \\
            }
        }
    }
    & & \\ \\

    \theutterance \stepcounter{utterance}  
    & & \multicolumn{4}{p{0.6\linewidth}}{
        \cellcolor[rgb]{0.9,0.9,0.9}{
            \makecell[{{p{\linewidth}}}]{
                \texttt{\tiny{[P1$\langle$GM]}}
                \texttt{Your message has been relayed to the other player.} \\
\\ 
\\ 
\texttt{You are currently playing round 5 of maximum 20.} \\
\texttt{You have collectively accumulated 1 of 12 penalties.} \\
\texttt{The other player sent this message:} \\
\\ 
\\ 
\texttt{"My 'C' is currently at (2,4). Considering the grid and the current positions, I think placing the 'C's at column 2 could be a good idea, since both our 'C's are already at that column. We could aim to place them at the same row, perhaps at a row that is relatively empty or close to the center. What do you think about placing the 'C's at row 4 or 5?"} \\
\\ 
\\ 
\texttt{What is your next command?} \\
            }
        }
    }
    & \\ \\

    \theutterance \stepcounter{utterance}  
    & \multicolumn{4}{p{0.6\linewidth}}{
        \cellcolor[rgb]{0.9,0.9,0.9}{
            \makecell[{{p{\linewidth}}}]{
                \texttt{\tiny{[P1$\rangle$GM]}}
                \texttt{SAY: Placing the 'C's at column 2 is a good idea, given their current positions. Between rows 4 and 5, I think row 5 might be a better choice. My grid has an 'I' at (4,2), a 'W' at (1,3), and an 'H' at (4,3), so row 5 seems relatively empty compared to row 4. If we place the 'C's at row 5, we can avoid potential conflicts with other objects. What do you think about placing the 'C's at (2,5)? Should we move our 'C's to this position and then discuss the placement of the 'H's?} \\
            }
        }
    }
    & & \\ \\

    \theutterance \stepcounter{utterance}  
    & & & \multicolumn{2}{p{0.3\linewidth}}{
        \cellcolor[rgb]{0.9,0.9,0.9}{
            \makecell[{{p{\linewidth}}}]{
                \texttt{\tiny{[GM$|$GM]}}
                \texttt{SAY: Placing the 'C's at column 2 is a good idea, given their current positions. Between rows 4 and 5, I think row 5 might be a better choice. My grid has an 'I' at (4,2), a 'W' at (1,3), and an 'H' at (4,3), so row 5 seems relatively empty compared to row 4. If we place the 'C's at row 5, we can avoid potential conflicts with other objects. What do you think about placing the 'C's at (2,5)? Should we move our 'C's to this position and then discuss the placement of the 'H's?} \\
            }
        }
    }
    & & \\ \\

    \theutterance \stepcounter{utterance}  
    & & \multicolumn{4}{p{0.6\linewidth}}{
        \cellcolor[rgb]{0.9,0.9,0.9}{
            \makecell[{{p{\linewidth}}}]{
                \texttt{\tiny{[GM$\rangle$P2]}}
                \texttt{Your message has been relayed to the other player.} \\
\\ 
\\ 
\texttt{You are currently playing round 6 of maximum 20.} \\
\texttt{You have collectively accumulated 1 of 12 penalties.} \\
\texttt{The other player sent this message:} \\
\\ 
\\ 
\texttt{"Placing the 'C's at column 2 is a good idea, given their current positions. Between rows 4 and 5, I think row 5 might be a better choice. My grid has an 'I' at (4,2), a 'W' at (1,3), and an 'H' at (4,3), so row 5 seems relatively empty compared to row 4. If we place the 'C's at row 5, we can avoid potential conflicts with other objects. What do you think about placing the 'C's at (2,5)? Should we move our 'C's to this position and then discuss the placement of the 'H's?"} \\
\\ 
\\ 
\texttt{What is your next command?} \\
            }
        }
    }
    & \\ \\

    \theutterance \stepcounter{utterance}  
    & & & \multicolumn{4}{p{0.6\linewidth}}{
        \cellcolor[rgb]{0.9,0.9,0.9}{
            \makecell[{{p{\linewidth}}}]{
                \texttt{\tiny{[GM$\langle$P2]}}
                \texttt{SAY: I agree that row 5 seems like a better choice, and placing the 'C's at (2,5) sounds like a good plan. I'll move my 'C' to (2,5). After that, we can discuss the 'H's. Since we're following alphabetical order, it makes sense to prioritize the 'H's next. What's your thoughts on where we should place the 'H's, considering their current positions on both our grids?} \\
            }
        }
    }
     \\ \\

    \theutterance \stepcounter{utterance}  
    & & & \multicolumn{2}{p{0.3\linewidth}}{
        \cellcolor[rgb]{0.9,0.9,0.9}{
            \makecell[{{p{\linewidth}}}]{
                \texttt{\tiny{[GM$|$GM]}}
                \texttt{SAY: I agree that row 5 seems like a better choice, and placing the 'C's at (2,5) sounds like a good plan. I'll move my 'C' to (2,5). After that, we can discuss the 'H's. Since we're following alphabetical order, it makes sense to prioritize the 'H's next. What's your thoughts on where we should place the 'H's, considering their current positions on both our grids?} \\
            }
        }
    }
    & & \\ \\

    \theutterance \stepcounter{utterance}  
    & & \multicolumn{4}{p{0.6\linewidth}}{
        \cellcolor[rgb]{0.9,0.9,0.9}{
            \makecell[{{p{\linewidth}}}]{
                \texttt{\tiny{[P1$\langle$GM]}}
                \texttt{Your message has been relayed to the other player.} \\
\\ 
\\ 
\texttt{You are currently playing round 6 of maximum 20.} \\
\texttt{You have collectively accumulated 1 of 12 penalties.} \\
\texttt{The other player sent this message:} \\
\\ 
\\ 
\texttt{"I agree that row 5 seems like a better choice, and placing the 'C's at (2,5) sounds like a good plan. I'll move my 'C' to (2,5). After that, we can discuss the 'H's. Since we're following alphabetical order, it makes sense to prioritize the 'H's next. What's your thoughts on where we should place the 'H's, considering their current positions on both our grids?"} \\
\\ 
\\ 
\texttt{What is your next command?} \\
            }
        }
    }
    & \\ \\

    \theutterance \stepcounter{utterance}  
    & \multicolumn{4}{p{0.6\linewidth}}{
        \cellcolor[rgb]{0.9,0.9,0.9}{
            \makecell[{{p{\linewidth}}}]{
                \texttt{\tiny{[P1$\rangle$GM]}}
                \texttt{SAY: Before we discuss the 'H's, I'll wait for your confirmation that your 'C' has been moved to (2,5). Once that's done, we can focus on the 'H's. My 'H' is currently at (4,3), and you mentioned earlier that your 'H' is at (1,6). To find a suitable position for the 'H's, let's consider the available space on our grids. Perhaps we could aim to place the 'H's at a row that is relatively empty, similar to what we did with the 'C's. What do you think about placing the 'H's at row 6, since your 'H' is already close to that row?} \\
            }
        }
    }
    & & \\ \\

    \theutterance \stepcounter{utterance}  
    & & & \multicolumn{2}{p{0.3\linewidth}}{
        \cellcolor[rgb]{0.9,0.9,0.9}{
            \makecell[{{p{\linewidth}}}]{
                \texttt{\tiny{[GM$|$GM]}}
                \texttt{SAY: Before we discuss the 'H's, I'll wait for your confirmation that your 'C' has been moved to (2,5). Once that's done, we can focus on the 'H's. My 'H' is currently at (4,3), and you mentioned earlier that your 'H' is at (1,6). To find a suitable position for the 'H's, let's consider the available space on our grids. Perhaps we could aim to place the 'H's at a row that is relatively empty, similar to what we did with the 'C's. What do you think about placing the 'H's at row 6, since your 'H' is already close to that row?} \\
            }
        }
    }
    & & \\ \\

    \theutterance \stepcounter{utterance}  
    & & \multicolumn{4}{p{0.6\linewidth}}{
        \cellcolor[rgb]{0.9,0.9,0.9}{
            \makecell[{{p{\linewidth}}}]{
                \texttt{\tiny{[GM$\rangle$P2]}}
                \texttt{Your message has been relayed to the other player.} \\
\\ 
\\ 
\texttt{You are currently playing round 7 of maximum 20.} \\
\texttt{You have collectively accumulated 1 of 12 penalties.} \\
\texttt{The other player sent this message:} \\
\\ 
\\ 
\texttt{"Before we discuss the 'H's, I'll wait for your confirmation that your 'C' has been moved to (2,5). Once that's done, we can focus on the 'H's. My 'H' is currently at (4,3), and you mentioned earlier that your 'H' is at (1,6). To find a suitable position for the 'H's, let's consider the available space on our grids. Perhaps we could aim to place the 'H's at a row that is relatively empty, similar to what we did with the 'C's. What do you think about placing the 'H's at row 6, since your 'H' is already close to that row?"} \\
\\ 
\\ 
\texttt{What is your next command?} \\
            }
        }
    }
    & \\ \\

    \theutterance \stepcounter{utterance}  
    & & & \multicolumn{4}{p{0.6\linewidth}}{
        \cellcolor[rgb]{0.9,0.9,0.9}{
            \makecell[{{p{\linewidth}}}]{
                \texttt{\tiny{[GM$\langle$P2]}}
                \texttt{SAY: I will move my 'C' to (2,5) now. MOVE: C, (2,5)} \\
            }
        }
    }
     \\ \\

    \theutterance \stepcounter{utterance}  
    & & & \multicolumn{2}{p{0.3\linewidth}}{
        \cellcolor[rgb]{0.9,0.9,0.9}{
            \makecell[{{p{\linewidth}}}]{
                \texttt{\tiny{[GM$|$GM]}}
                \texttt{SAY: I will move my 'C' to (2,5) now. MOVE: C, (2,5)} \\
            }
        }
    }
    & & \\ \\

    \theutterance \stepcounter{utterance}  
    & & & \multicolumn{2}{p{0.3\linewidth}}{
        \cellcolor[rgb]{0.9,0.9,0.9}{
            \makecell[{{p{\linewidth}}}]{
                \texttt{\tiny{[GM$|$GM]}}
                \texttt{Invalid response format: SAY: I will move my 'C' to (2,5) now. MOVE: C, (2,5)} \\
            }
        }
    }
    & & \\ \\

    \theutterance \stepcounter{utterance}  
    & & \multicolumn{4}{p{0.6\linewidth}}{
        \cellcolor[rgb]{0.9,0.9,0.9}{
            \makecell[{{p{\linewidth}}}]{
                \texttt{\tiny{[GM$\rangle$P2]}}
                \texttt{Penalty: Your message contains more than one command!} \\
\\ 
\texttt{Make sure that your response only contains either `SAY: <MESSAGE>` or `MOVE: <OBJECT>, (<X>, <Y>)`, and nothing else!} \\
\\ 
\texttt{You have collectively accumulated 2 of 12 penalties.} \\
\texttt{Please try again!} \\
            }
        }
    }
    & \\ \\

    \theutterance \stepcounter{utterance}  
    & & & \multicolumn{4}{p{0.6\linewidth}}{
        \cellcolor[rgb]{0.9,0.9,0.9}{
            \makecell[{{p{\linewidth}}}]{
                \texttt{\tiny{[GM$\langle$P2]}}
                \texttt{MOVE: C, (2,5)} \\
            }
        }
    }
     \\ \\

    \theutterance \stepcounter{utterance}  
    & & & \multicolumn{2}{p{0.3\linewidth}}{
        \cellcolor[rgb]{0.9,0.9,0.9}{
            \makecell[{{p{\linewidth}}}]{
                \texttt{\tiny{[GM$|$GM]}}
                \texttt{MOVE: C, (2,5)} \\
            }
        }
    }
    & & \\ \\

    \theutterance \stepcounter{utterance}  
    & & & \multicolumn{2}{p{0.3\linewidth}}{
        \cellcolor[rgb]{0.9,0.9,0.9}{
            \makecell[{{p{\linewidth}}}]{
                \texttt{\tiny{[GM$|$GM]}}
                \texttt{Penalty: (2,5) is not empty, but contains '┴'.} \\
\\ 
\texttt{You have collectively accumulated 3 of 12 penalties.} \\
\texttt{Please try again!} \\
            }
        }
    }
    & & \\ \\

    \theutterance \stepcounter{utterance}  
    & & \multicolumn{4}{p{0.6\linewidth}}{
        \cellcolor[rgb]{0.9,0.9,0.9}{
            \makecell[{{p{\linewidth}}}]{
                \texttt{\tiny{[GM$\rangle$P2]}}
                \texttt{Penalty: (2,5) is not empty, but contains '┴'.} \\
\\ 
\texttt{You have collectively accumulated 3 of 12 penalties.} \\
\texttt{Please try again!} \\
            }
        }
    }
    & \\ \\

    \theutterance \stepcounter{utterance}  
    & & & \multicolumn{4}{p{0.6\linewidth}}{
        \cellcolor[rgb]{0.9,0.9,0.9}{
            \makecell[{{p{\linewidth}}}]{
                \texttt{\tiny{[GM$\langle$P2]}}
                \texttt{MOVE: C, (2,5) is not a valid move. Let me try again. MOVE: C, (2,4) is also not needed, since it is already there. I will move it to a different position to make space. MOVE: C, (3,4)} \\
            }
        }
    }
     \\ \\

    \theutterance \stepcounter{utterance}  
    & & & \multicolumn{2}{p{0.3\linewidth}}{
        \cellcolor[rgb]{0.9,0.9,0.9}{
            \makecell[{{p{\linewidth}}}]{
                \texttt{\tiny{[GM$|$GM]}}
                \texttt{MOVE: C, (2,5) is not a valid move. Let me try again. MOVE: C, (2,4) is also not needed, since it is already there. I will move it to a different position to make space. MOVE: C, (3,4)} \\
            }
        }
    }
    & & \\ \\

    \theutterance \stepcounter{utterance}  
    & & & \multicolumn{2}{p{0.3\linewidth}}{
        \cellcolor[rgb]{0.9,0.9,0.9}{
            \makecell[{{p{\linewidth}}}]{
                \texttt{\tiny{[GM$|$GM]}}
                \texttt{Invalid format: head is not empty: 'MOVE: C, (2,5) is not a valid move. Let me try again. MOVE: C, (2,4) is also not needed, since it is already there. I will move it to a different position to make space. '} \\
            }
        }
    }
    & & \\ \\

    \theutterance \stepcounter{utterance}  
    & & \multicolumn{4}{p{0.6\linewidth}}{
        \cellcolor[rgb]{0.9,0.9,0.9}{
            \makecell[{{p{\linewidth}}}]{
                \texttt{\tiny{[GM$\rangle$P2]}}
                \texttt{Penalty: Your message must not contain anything before the command!} \\
\\ 
\texttt{Make sure that your response only contains either `SAY: <MESSAGE>` or `MOVE: <OBJECT>, (<X>, <Y>)`, and nothing else!} \\
\\ 
\texttt{You have collectively accumulated 4 of 12 penalties.} \\
\texttt{Please try again!} \\
            }
        }
    }
    & \\ \\

    \theutterance \stepcounter{utterance}  
    & & & \multicolumn{4}{p{0.6\linewidth}}{
        \cellcolor[rgb]{0.9,0.9,0.9}{
            \makecell[{{p{\linewidth}}}]{
                \texttt{\tiny{[GM$\langle$P2]}}
                \texttt{MOVE: C, (2,5)} \\
            }
        }
    }
     \\ \\

    \theutterance \stepcounter{utterance}  
    & & & \multicolumn{2}{p{0.3\linewidth}}{
        \cellcolor[rgb]{0.9,0.9,0.9}{
            \makecell[{{p{\linewidth}}}]{
                \texttt{\tiny{[GM$|$GM]}}
                \texttt{MOVE: C, (2,5)} \\
            }
        }
    }
    & & \\ \\

    \theutterance \stepcounter{utterance}  
    & & & \multicolumn{2}{p{0.3\linewidth}}{
        \cellcolor[rgb]{0.9,0.9,0.9}{
            \makecell[{{p{\linewidth}}}]{
                \texttt{\tiny{[GM$|$GM]}}
                \texttt{Penalty: (2,5) is not empty, but contains '┴'.} \\
\\ 
\texttt{You have collectively accumulated 5 of 12 penalties.} \\
\texttt{Please try again!} \\
            }
        }
    }
    & & \\ \\

    \theutterance \stepcounter{utterance}  
    & & \multicolumn{4}{p{0.6\linewidth}}{
        \cellcolor[rgb]{0.9,0.9,0.9}{
            \makecell[{{p{\linewidth}}}]{
                \texttt{\tiny{[GM$\rangle$P2]}}
                \texttt{Penalty: (2,5) is not empty, but contains '┴'.} \\
\\ 
\texttt{You have collectively accumulated 5 of 12 penalties.} \\
\texttt{Please try again!} \\
            }
        }
    }
    & \\ \\

    \theutterance \stepcounter{utterance}  
    & & & \multicolumn{4}{p{0.6\linewidth}}{
        \cellcolor[rgb]{0.9,0.9,0.9}{
            \makecell[{{p{\linewidth}}}]{
                \texttt{\tiny{[GM$\langle$P2]}}
                \texttt{MOVE: C, (4,4)} \\
            }
        }
    }
     \\ \\

    \theutterance \stepcounter{utterance}  
    & & & \multicolumn{2}{p{0.3\linewidth}}{
        \cellcolor[rgb]{0.9,0.9,0.9}{
            \makecell[{{p{\linewidth}}}]{
                \texttt{\tiny{[GM$|$GM]}}
                \texttt{MOVE: C, (4,4)} \\
            }
        }
    }
    & & \\ \\

    \theutterance \stepcounter{utterance}  
    & & & \multicolumn{2}{p{0.3\linewidth}}{
        \cellcolor[rgb]{0.9,0.9,0.9}{
            \makecell[{{p{\linewidth}}}]{
                \texttt{\tiny{[GM$|$GM]}}
                \texttt{Penalty: (4,4) is not empty, but contains '├'.} \\
\\ 
\texttt{You have collectively accumulated 6 of 12 penalties.} \\
\texttt{Please try again!} \\
            }
        }
    }
    & & \\ \\

    \theutterance \stepcounter{utterance}  
    & & \multicolumn{4}{p{0.6\linewidth}}{
        \cellcolor[rgb]{0.9,0.9,0.9}{
            \makecell[{{p{\linewidth}}}]{
                \texttt{\tiny{[GM$\rangle$P2]}}
                \texttt{Penalty: (4,4) is not empty, but contains '├'.} \\
\\ 
\texttt{You have collectively accumulated 6 of 12 penalties.} \\
\texttt{Please try again!} \\
            }
        }
    }
    & \\ \\

    \theutterance \stepcounter{utterance}  
    & & & \multicolumn{4}{p{0.6\linewidth}}{
        \cellcolor[rgb]{0.9,0.9,0.9}{
            \makecell[{{p{\linewidth}}}]{
                \texttt{\tiny{[GM$\langle$P2]}}
                \texttt{MOVE: C, (5,4)} \\
            }
        }
    }
     \\ \\

    \theutterance \stepcounter{utterance}  
    & & & \multicolumn{2}{p{0.3\linewidth}}{
        \cellcolor[rgb]{0.9,0.9,0.9}{
            \makecell[{{p{\linewidth}}}]{
                \texttt{\tiny{[GM$|$GM]}}
                \texttt{MOVE: C, (5,4)} \\
            }
        }
    }
    & & \\ \\

    \theutterance \stepcounter{utterance}  
    & & & \multicolumn{2}{p{0.3\linewidth}}{
        \cellcolor[rgb]{0.9,0.9,0.9}{
            \makecell[{{p{\linewidth}}}]{
                \texttt{\tiny{[GM$|$GM]}}
                \texttt{Penalty: (5,4) is not empty, but contains '─'.} \\
\\ 
\texttt{You have collectively accumulated 7 of 12 penalties.} \\
\texttt{Please try again!} \\
            }
        }
    }
    & & \\ \\

    \theutterance \stepcounter{utterance}  
    & & \multicolumn{4}{p{0.6\linewidth}}{
        \cellcolor[rgb]{0.9,0.9,0.9}{
            \makecell[{{p{\linewidth}}}]{
                \texttt{\tiny{[GM$\rangle$P2]}}
                \texttt{Penalty: (5,4) is not empty, but contains '─'.} \\
\\ 
\texttt{You have collectively accumulated 7 of 12 penalties.} \\
\texttt{Please try again!} \\
            }
        }
    }
    & \\ \\

    \theutterance \stepcounter{utterance}  
    & & & \multicolumn{4}{p{0.6\linewidth}}{
        \cellcolor[rgb]{0.9,0.9,0.9}{
            \makecell[{{p{\linewidth}}}]{
                \texttt{\tiny{[GM$\langle$P2]}}
                \texttt{MOVE: C, (3,5)} \\
            }
        }
    }
     \\ \\

    \theutterance \stepcounter{utterance}  
    & & & \multicolumn{2}{p{0.3\linewidth}}{
        \cellcolor[rgb]{0.9,0.9,0.9}{
            \makecell[{{p{\linewidth}}}]{
                \texttt{\tiny{[GM$|$GM]}}
                \texttt{MOVE: C, (3,5)} \\
            }
        }
    }
    & & \\ \\

    \theutterance \stepcounter{utterance}  
    & & & \multicolumn{2}{p{0.3\linewidth}}{
        \cellcolor[rgb]{0.9,0.9,0.9}{
            \makecell[{{p{\linewidth}}}]{
                \texttt{\tiny{[GM$|$GM]}}
                \texttt{Penalty: (3,5) is not empty, but contains '─'.} \\
\\ 
\texttt{You have collectively accumulated 8 of 12 penalties.} \\
\texttt{Please try again!} \\
            }
        }
    }
    & & \\ \\

    \theutterance \stepcounter{utterance}  
    & & \multicolumn{4}{p{0.6\linewidth}}{
        \cellcolor[rgb]{0.9,0.9,0.9}{
            \makecell[{{p{\linewidth}}}]{
                \texttt{\tiny{[GM$\rangle$P2]}}
                \texttt{Penalty: (3,5) is not empty, but contains '─'.} \\
\\ 
\texttt{You have collectively accumulated 8 of 12 penalties.} \\
\texttt{Please try again!} \\
            }
        }
    }
    & \\ \\

    \theutterance \stepcounter{utterance}  
    & & & \multicolumn{4}{p{0.6\linewidth}}{
        \cellcolor[rgb]{0.9,0.9,0.9}{
            \makecell[{{p{\linewidth}}}]{
                \texttt{\tiny{[GM$\langle$P2]}}
                \texttt{MOVE: C, (5,5)} \\
            }
        }
    }
     \\ \\

    \theutterance \stepcounter{utterance}  
    & & & \multicolumn{2}{p{0.3\linewidth}}{
        \cellcolor[rgb]{0.9,0.9,0.9}{
            \makecell[{{p{\linewidth}}}]{
                \texttt{\tiny{[GM$|$GM]}}
                \texttt{MOVE: C, (5,5)} \\
            }
        }
    }
    & & \\ \\

    \theutterance \stepcounter{utterance}  
    & & & \multicolumn{2}{p{0.3\linewidth}}{
        \cellcolor[rgb]{0.9,0.9,0.9}{
            \makecell[{{p{\linewidth}}}]{
                \texttt{\tiny{[GM$|$GM]}}
                \texttt{Penalty: (5,5) is not empty, but contains 'I'.} \\
\\ 
\texttt{You have collectively accumulated 9 of 12 penalties.} \\
\texttt{Please try again!} \\
            }
        }
    }
    & & \\ \\

    \theutterance \stepcounter{utterance}  
    & & \multicolumn{4}{p{0.6\linewidth}}{
        \cellcolor[rgb]{0.9,0.9,0.9}{
            \makecell[{{p{\linewidth}}}]{
                \texttt{\tiny{[GM$\rangle$P2]}}
                \texttt{Penalty: (5,5) is not empty, but contains 'I'.} \\
\\ 
\texttt{You have collectively accumulated 9 of 12 penalties.} \\
\texttt{Please try again!} \\
            }
        }
    }
    & \\ \\

    \theutterance \stepcounter{utterance}  
    & & & \multicolumn{4}{p{0.6\linewidth}}{
        \cellcolor[rgb]{0.9,0.9,0.9}{
            \makecell[{{p{\linewidth}}}]{
                \texttt{\tiny{[GM$\langle$P2]}}
                \texttt{MOVE: C, (6,5)} \\
            }
        }
    }
     \\ \\

    \theutterance \stepcounter{utterance}  
    & & & \multicolumn{2}{p{0.3\linewidth}}{
        \cellcolor[rgb]{0.9,0.9,0.9}{
            \makecell[{{p{\linewidth}}}]{
                \texttt{\tiny{[GM$|$GM]}}
                \texttt{MOVE: C, (6,5)} \\
            }
        }
    }
    & & \\ \\

    \theutterance \stepcounter{utterance}  
    & & & \multicolumn{2}{p{0.3\linewidth}}{
        \cellcolor[rgb]{0.9,0.9,0.9}{
            \makecell[{{p{\linewidth}}}]{
                \texttt{\tiny{[GM$|$GM]}}
                \texttt{C moved to (6, 5)} \\
            }
        }
    }
    & & \\ \\

    \theutterance \stepcounter{utterance}  
    & & \multicolumn{4}{p{0.6\linewidth}}{
        \cellcolor[rgb]{0.9,0.9,0.9}{
            \makecell[{{p{\linewidth}}}]{
                \texttt{\tiny{[P1$\langle$GM]}}
                \texttt{Your message has been relayed to the other player.} \\
\\ 
\\ 
\texttt{You are currently playing round 7 of maximum 20.} \\
\texttt{You have collectively accumulated 9 of 12 penalties.} \\
\texttt{The other player moved an object on their grid.} \\
\\ 
\\ 
\texttt{What is your next command?} \\
            }
        }
    }
    & \\ \\

    \theutterance \stepcounter{utterance}  
    & \multicolumn{4}{p{0.6\linewidth}}{
        \cellcolor[rgb]{0.9,0.9,0.9}{
            \makecell[{{p{\linewidth}}}]{
                \texttt{\tiny{[P1$\rangle$GM]}}
                \texttt{SAY: I've noticed that the penalty count has increased significantly. To avoid further penalties, let's ensure that we're making valid moves and communicating clearly. Regarding the 'H's, I suggested placing them at row 6. Could you please confirm if this is a viable option for your grid, considering the new position of your 'C' and any other objects that may be in the way? Additionally, I'll review my previous moves to ensure I'm not missing any important details. Please share your thoughts on the 'H's placement, and we can proceed from there.} \\
            }
        }
    }
    & & \\ \\

    \theutterance \stepcounter{utterance}  
    & & & \multicolumn{2}{p{0.3\linewidth}}{
        \cellcolor[rgb]{0.9,0.9,0.9}{
            \makecell[{{p{\linewidth}}}]{
                \texttt{\tiny{[GM$|$GM]}}
                \texttt{SAY: I've noticed that the penalty count has increased significantly. To avoid further penalties, let's ensure that we're making valid moves and communicating clearly. Regarding the 'H's, I suggested placing them at row 6. Could you please confirm if this is a viable option for your grid, considering the new position of your 'C' and any other objects that may be in the way? Additionally, I'll review my previous moves to ensure I'm not missing any important details. Please share your thoughts on the 'H's placement, and we can proceed from there.} \\
            }
        }
    }
    & & \\ \\

    \theutterance \stepcounter{utterance}  
    & & \multicolumn{4}{p{0.6\linewidth}}{
        \cellcolor[rgb]{0.9,0.9,0.9}{
            \makecell[{{p{\linewidth}}}]{
                \texttt{\tiny{[GM$\rangle$P2]}}
                \texttt{Moved 'C' to (6,5) successfully. Your updated grid looks like this:} \\
\texttt{```} \\
\texttt{ 1234567} \\
\texttt{╔═╤═╤═╤═╗} \\
\texttt{║◌│◌│◌│◌║ 1} \\
\texttt{╟─┼─┤◌│◌║ 2} \\
\texttt{║◌│W│◌│◌║ 3} \\
\texttt{║◌│◌├─┴─╢ 4} \\
\texttt{╟─┴─┤IC◌║ 5} \\
\texttt{║◌H◌└─┬─╢ 6} \\
\texttt{║◌◌T◌◌│◌║ 7} \\
\texttt{╚═════╧═╝} \\
\texttt{```} \\
\texttt{You are currently playing round 8 of maximum 20.} \\
\texttt{You have collectively accumulated 9 of 12 penalties.} \\
\texttt{The other player sent this message:} \\
\\ 
\\ 
\texttt{"I've noticed that the penalty count has increased significantly. To avoid further penalties, let's ensure that we're making valid moves and communicating clearly. Regarding the 'H's, I suggested placing them at row 6. Could you please confirm if this is a viable option for your grid, considering the new position of your 'C' and any other objects that may be in the way? Additionally, I'll review my previous moves to ensure I'm not missing any important details. Please share your thoughts on the 'H's placement, and we can proceed from there."} \\
\\ 
\\ 
\texttt{What is your next command?} \\
            }
        }
    }
    & \\ \\

    \theutterance \stepcounter{utterance}  
    & & & \multicolumn{4}{p{0.6\linewidth}}{
        \cellcolor[rgb]{0.9,0.9,0.9}{
            \makecell[{{p{\linewidth}}}]{
                \texttt{\tiny{[GM$\langle$P2]}}
                \texttt{SAY: Now that my 'C' is at (6,5), I can focus on the 'H's. Placing the 'H's at row 6 seems like a good idea. My 'H' is currently at (1,6), so it's already at the desired row. I can try to move it to a column that works for both our grids. What do you think about placing the 'H's at column 1, since my 'H' is already there?} \\
            }
        }
    }
     \\ \\

    \theutterance \stepcounter{utterance}  
    & & & \multicolumn{2}{p{0.3\linewidth}}{
        \cellcolor[rgb]{0.9,0.9,0.9}{
            \makecell[{{p{\linewidth}}}]{
                \texttt{\tiny{[GM$|$GM]}}
                \texttt{SAY: Now that my 'C' is at (6,5), I can focus on the 'H's. Placing the 'H's at row 6 seems like a good idea. My 'H' is currently at (1,6), so it's already at the desired row. I can try to move it to a column that works for both our grids. What do you think about placing the 'H's at column 1, since my 'H' is already there?} \\
            }
        }
    }
    & & \\ \\

    \theutterance \stepcounter{utterance}  
    & & \multicolumn{4}{p{0.6\linewidth}}{
        \cellcolor[rgb]{0.9,0.9,0.9}{
            \makecell[{{p{\linewidth}}}]{
                \texttt{\tiny{[P1$\langle$GM]}}
                \texttt{Your message has been relayed to the other player.} \\
\\ 
\\ 
\texttt{You are currently playing round 8 of maximum 20.} \\
\texttt{You have collectively accumulated 9 of 12 penalties.} \\
\texttt{The other player sent this message:} \\
\\ 
\\ 
\texttt{"Now that my 'C' is at (6,5), I can focus on the 'H's. Placing the 'H's at row 6 seems like a good idea. My 'H' is currently at (1,6), so it's already at the desired row. I can try to move it to a column that works for both our grids. What do you think about placing the 'H's at column 1, since my 'H' is already there?"} \\
\\ 
\\ 
\texttt{What is your next command?} \\
            }
        }
    }
    & \\ \\

    \theutterance \stepcounter{utterance}  
    & \multicolumn{4}{p{0.6\linewidth}}{
        \cellcolor[rgb]{0.9,0.9,0.9}{
            \makecell[{{p{\linewidth}}}]{
                \texttt{\tiny{[P1$\rangle$GM]}}
                \texttt{SAY: I think there might be some confusion. You mentioned your 'C' is at (6,5), but earlier we discussed placing the 'C's at (2,5). Could you please clarify the correct position of your 'C'? Additionally, regarding the 'H's, placing them at column 1 could work, but I need to consider the position of my 'H' at (4,3). I'll need to move it to row 6 and column 1. Before I make any moves, I want to ensure we're on the same page. Can you please confirm the correct position of your 'C' and if you're ready for me to move my 'H' to (1,6)?} \\
            }
        }
    }
    & & \\ \\

    \theutterance \stepcounter{utterance}  
    & & & \multicolumn{2}{p{0.3\linewidth}}{
        \cellcolor[rgb]{0.9,0.9,0.9}{
            \makecell[{{p{\linewidth}}}]{
                \texttt{\tiny{[GM$|$GM]}}
                \texttt{SAY: I think there might be some confusion. You mentioned your 'C' is at (6,5), but earlier we discussed placing the 'C's at (2,5). Could you please clarify the correct position of your 'C'? Additionally, regarding the 'H's, placing them at column 1 could work, but I need to consider the position of my 'H' at (4,3). I'll need to move it to row 6 and column 1. Before I make any moves, I want to ensure we're on the same page. Can you please confirm the correct position of your 'C' and if you're ready for me to move my 'H' to (1,6)?} \\
            }
        }
    }
    & & \\ \\

    \theutterance \stepcounter{utterance}  
    & & \multicolumn{4}{p{0.6\linewidth}}{
        \cellcolor[rgb]{0.9,0.9,0.9}{
            \makecell[{{p{\linewidth}}}]{
                \texttt{\tiny{[GM$\rangle$P2]}}
                \texttt{Your message has been relayed to the other player.} \\
\\ 
\\ 
\texttt{You are currently playing round 9 of maximum 20.} \\
\texttt{You have collectively accumulated 9 of 12 penalties.} \\
\texttt{The other player sent this message:} \\
\\ 
\\ 
\texttt{"I think there might be some confusion. You mentioned your 'C' is at (6,5), but earlier we discussed placing the 'C's at (2,5). Could you please clarify the correct position of your 'C'? Additionally, regarding the 'H's, placing them at column 1 could work, but I need to consider the position of my 'H' at (4,3). I'll need to move it to row 6 and column 1. Before I make any moves, I want to ensure we're on the same page. Can you please confirm the correct position of your 'C' and if you're ready for me to move my 'H' to (1,6)?"} \\
\\ 
\\ 
\texttt{What is your next command?} \\
            }
        }
    }
    & \\ \\

    \theutterance \stepcounter{utterance}  
    & & & \multicolumn{4}{p{0.6\linewidth}}{
        \cellcolor[rgb]{0.9,0.9,0.9}{
            \makecell[{{p{\linewidth}}}]{
                \texttt{\tiny{[GM$\langle$P2]}}
                \texttt{SAY: You are correct, we initially discussed placing the 'C's at (2,5), but due to the penalties and moves, my 'C' is now actually at (6,5). To get back on track, I can try to move my 'C' to the originally planned position (2,5), but I need to check if the path is clear. As for the 'H's, yes, I'm ready for you to move your 'H' to (1,6), which would put both 'H's in the same column. Please proceed with moving your 'H'.} \\
            }
        }
    }
     \\ \\

    \theutterance \stepcounter{utterance}  
    & & & \multicolumn{2}{p{0.3\linewidth}}{
        \cellcolor[rgb]{0.9,0.9,0.9}{
            \makecell[{{p{\linewidth}}}]{
                \texttt{\tiny{[GM$|$GM]}}
                \texttt{SAY: You are correct, we initially discussed placing the 'C's at (2,5), but due to the penalties and moves, my 'C' is now actually at (6,5). To get back on track, I can try to move my 'C' to the originally planned position (2,5), but I need to check if the path is clear. As for the 'H's, yes, I'm ready for you to move your 'H' to (1,6), which would put both 'H's in the same column. Please proceed with moving your 'H'.} \\
            }
        }
    }
    & & \\ \\

    \theutterance \stepcounter{utterance}  
    & & \multicolumn{4}{p{0.6\linewidth}}{
        \cellcolor[rgb]{0.9,0.9,0.9}{
            \makecell[{{p{\linewidth}}}]{
                \texttt{\tiny{[P1$\langle$GM]}}
                \texttt{Your message has been relayed to the other player.} \\
\\ 
\\ 
\texttt{You are currently playing round 9 of maximum 20.} \\
\texttt{You have collectively accumulated 9 of 12 penalties.} \\
\texttt{The other player sent this message:} \\
\\ 
\\ 
\texttt{"You are correct, we initially discussed placing the 'C's at (2,5), but due to the penalties and moves, my 'C' is now actually at (6,5). To get back on track, I can try to move my 'C' to the originally planned position (2,5), but I need to check if the path is clear. As for the 'H's, yes, I'm ready for you to move your 'H' to (1,6), which would put both 'H's in the same column. Please proceed with moving your 'H'."} \\
\\ 
\\ 
\texttt{What is your next command?} \\
            }
        }
    }
    & \\ \\

    \theutterance \stepcounter{utterance}  
    & \multicolumn{4}{p{0.6\linewidth}}{
        \cellcolor[rgb]{0.9,0.9,0.9}{
            \makecell[{{p{\linewidth}}}]{
                \texttt{\tiny{[P1$\rangle$GM]}}
                \texttt{MOVE: H, (1, 6)} \\
            }
        }
    }
    & & \\ \\

    \theutterance \stepcounter{utterance}  
    & & & \multicolumn{2}{p{0.3\linewidth}}{
        \cellcolor[rgb]{0.9,0.9,0.9}{
            \makecell[{{p{\linewidth}}}]{
                \texttt{\tiny{[GM$|$GM]}}
                \texttt{MOVE: H, (1, 6)} \\
            }
        }
    }
    & & \\ \\

    \theutterance \stepcounter{utterance}  
    & & & \multicolumn{2}{p{0.3\linewidth}}{
        \cellcolor[rgb]{0.9,0.9,0.9}{
            \makecell[{{p{\linewidth}}}]{
                \texttt{\tiny{[GM$|$GM]}}
                \texttt{H moved to (1, 6)} \\
            }
        }
    }
    & & \\ \\

    \theutterance \stepcounter{utterance}  
    & & \multicolumn{4}{p{0.6\linewidth}}{
        \cellcolor[rgb]{0.9,0.9,0.9}{
            \makecell[{{p{\linewidth}}}]{
                \texttt{\tiny{[GM$\rangle$P2]}}
                \texttt{Your message has been relayed to the other player.} \\
\\ 
\\ 
\texttt{You are currently playing round 10 of maximum 20.} \\
\texttt{You have collectively accumulated 9 of 12 penalties.} \\
\texttt{The other player moved an object on their grid.} \\
\\ 
\\ 
\texttt{What is your next command?} \\
            }
        }
    }
    & \\ \\

    \theutterance \stepcounter{utterance}  
    & & & \multicolumn{4}{p{0.6\linewidth}}{
        \cellcolor[rgb]{0.9,0.9,0.9}{
            \makecell[{{p{\linewidth}}}]{
                \texttt{\tiny{[GM$\langle$P2]}}
                \texttt{SAY: Now that you've moved your object, can you confirm if your 'H' is indeed at (1,6)? If so, we can proceed with discussing the next steps for the remaining objects, such as the 'I's and 'W's, to get them into a consistent placement.} \\
            }
        }
    }
     \\ \\

    \theutterance \stepcounter{utterance}  
    & & & \multicolumn{2}{p{0.3\linewidth}}{
        \cellcolor[rgb]{0.9,0.9,0.9}{
            \makecell[{{p{\linewidth}}}]{
                \texttt{\tiny{[GM$|$GM]}}
                \texttt{SAY: Now that you've moved your object, can you confirm if your 'H' is indeed at (1,6)? If so, we can proceed with discussing the next steps for the remaining objects, such as the 'I's and 'W's, to get them into a consistent placement.} \\
            }
        }
    }
    & & \\ \\

    \theutterance \stepcounter{utterance}  
    & & \multicolumn{4}{p{0.6\linewidth}}{
        \cellcolor[rgb]{0.9,0.9,0.9}{
            \makecell[{{p{\linewidth}}}]{
                \texttt{\tiny{[P1$\langle$GM]}}
                \texttt{Moved 'H' to (1,6) successfully. Your updated grid looks like this:} \\
\texttt{```} \\
\texttt{ 1234567} \\
\texttt{╔═╤═╤═╤═╗} \\
\texttt{║◌│◌│◌│◌║ 1} \\
\texttt{╟─┼─┤I│◌║ 2} \\
\texttt{║W│◌│◌│◌║ 3} \\
\texttt{║◌│◌├─┴─╢ 4} \\
\texttt{╟─┴─┤◌◌◌║ 5} \\
\texttt{║HC◌└─┬─╢ 6} \\
\texttt{║◌◌T◌◌│◌║ 7} \\
\texttt{╚═════╧═╝} \\
\texttt{```} \\
\texttt{You are currently playing round 10 of maximum 20.} \\
\texttt{You have collectively accumulated 9 of 12 penalties.} \\
\texttt{The other player sent this message:} \\
\\ 
\\ 
\texttt{"Now that you've moved your object, can you confirm if your 'H' is indeed at (1,6)? If so, we can proceed with discussing the next steps for the remaining objects, such as the 'I's and 'W's, to get them into a consistent placement."} \\
\\ 
\\ 
\texttt{What is your next command?} \\
            }
        }
    }
    & \\ \\

    \theutterance \stepcounter{utterance}  
    & \multicolumn{4}{p{0.6\linewidth}}{
        \cellcolor[rgb]{0.9,0.9,0.9}{
            \makecell[{{p{\linewidth}}}]{
                \texttt{\tiny{[P1$\rangle$GM]}}
                \texttt{SAY: Yes, I confirm that my 'H' is now at (1,6). With the 'H's in the same column, let's focus on the 'I's. My 'I' is currently at (4,2), and you didn't mention the position of your 'I' earlier. Could you please share the current position of your 'I'? Once we have that information, we can discuss a suitable position for the 'I's. Perhaps we could aim to place them at a specific row or column that works for both our grids. After the 'I's, we can move on to the 'W's. What are your thoughts on how to proceed with the 'I's?} \\
            }
        }
    }
    & & \\ \\

    \theutterance \stepcounter{utterance}  
    & & & \multicolumn{2}{p{0.3\linewidth}}{
        \cellcolor[rgb]{0.9,0.9,0.9}{
            \makecell[{{p{\linewidth}}}]{
                \texttt{\tiny{[GM$|$GM]}}
                \texttt{SAY: Yes, I confirm that my 'H' is now at (1,6). With the 'H's in the same column, let's focus on the 'I's. My 'I' is currently at (4,2), and you didn't mention the position of your 'I' earlier. Could you please share the current position of your 'I'? Once we have that information, we can discuss a suitable position for the 'I's. Perhaps we could aim to place them at a specific row or column that works for both our grids. After the 'I's, we can move on to the 'W's. What are your thoughts on how to proceed with the 'I's?} \\
            }
        }
    }
    & & \\ \\

    \theutterance \stepcounter{utterance}  
    & & \multicolumn{4}{p{0.6\linewidth}}{
        \cellcolor[rgb]{0.9,0.9,0.9}{
            \makecell[{{p{\linewidth}}}]{
                \texttt{\tiny{[GM$\rangle$P2]}}
                \texttt{Your message has been relayed to the other player.} \\
\\ 
\\ 
\texttt{You are currently playing round 11 of maximum 20.} \\
\texttt{You have collectively accumulated 9 of 12 penalties.} \\
\texttt{The other player sent this message:} \\
\\ 
\\ 
\texttt{"Yes, I confirm that my 'H' is now at (1,6). With the 'H's in the same column, let's focus on the 'I's. My 'I' is currently at (4,2), and you didn't mention the position of your 'I' earlier. Could you please share the current position of your 'I'? Once we have that information, we can discuss a suitable position for the 'I's. Perhaps we could aim to place them at a specific row or column that works for both our grids. After the 'I's, we can move on to the 'W's. What are your thoughts on how to proceed with the 'I's?"} \\
\\ 
\\ 
\texttt{What is your next command?} \\
            }
        }
    }
    & \\ \\

    \theutterance \stepcounter{utterance}  
    & & & \multicolumn{4}{p{0.6\linewidth}}{
        \cellcolor[rgb]{0.9,0.9,0.9}{
            \makecell[{{p{\linewidth}}}]{
                \texttt{\tiny{[GM$\langle$P2]}}
                \texttt{SAY: My 'I' is currently at (1,5). Considering the positions of the 'H's and the 'I's, we could try to place the 'I's at a row that is relatively close to the 'H's, such as row 5 or 6, to maintain some consistency. Alternatively, we could aim to place the 'I's at a specific column, but I think row{-}based placement might be more feasible given the current grid layout. What do you think about placing the 'I's at row 5?} \\
            }
        }
    }
     \\ \\

    \theutterance \stepcounter{utterance}  
    & & & \multicolumn{2}{p{0.3\linewidth}}{
        \cellcolor[rgb]{0.9,0.9,0.9}{
            \makecell[{{p{\linewidth}}}]{
                \texttt{\tiny{[GM$|$GM]}}
                \texttt{SAY: My 'I' is currently at (1,5). Considering the positions of the 'H's and the 'I's, we could try to place the 'I's at a row that is relatively close to the 'H's, such as row 5 or 6, to maintain some consistency. Alternatively, we could aim to place the 'I's at a specific column, but I think row{-}based placement might be more feasible given the current grid layout. What do you think about placing the 'I's at row 5?} \\
            }
        }
    }
    & & \\ \\

    \theutterance \stepcounter{utterance}  
    & & \multicolumn{4}{p{0.6\linewidth}}{
        \cellcolor[rgb]{0.9,0.9,0.9}{
            \makecell[{{p{\linewidth}}}]{
                \texttt{\tiny{[P1$\langle$GM]}}
                \texttt{Your message has been relayed to the other player.} \\
\\ 
\\ 
\texttt{You are currently playing round 11 of maximum 20.} \\
\texttt{You have collectively accumulated 9 of 12 penalties.} \\
\texttt{The other player sent this message:} \\
\\ 
\\ 
\texttt{"My 'I' is currently at (1,5). Considering the positions of the 'H's and the 'I's, we could try to place the 'I's at a row that is relatively close to the 'H's, such as row 5 or 6, to maintain some consistency. Alternatively, we could aim to place the 'I's at a specific column, but I think row{-}based placement might be more feasible given the current grid layout. What do you think about placing the 'I's at row 5?"} \\
\\ 
\\ 
\texttt{What is your next command?} \\
            }
        }
    }
    & \\ \\

    \theutterance \stepcounter{utterance}  
    & \multicolumn{4}{p{0.6\linewidth}}{
        \cellcolor[rgb]{0.9,0.9,0.9}{
            \makecell[{{p{\linewidth}}}]{
                \texttt{\tiny{[P1$\rangle$GM]}}
                \texttt{SAY: Placing the 'I's at row 5 seems like a reasonable approach, considering the current positions of the 'H's and 'I's. Since my 'I' is at (4,2) and your 'I' is at (1,5), we would need to move my 'I' to row 5. To maintain consistency, let's aim to place the 'I's at the same column as well. How about we place the 'I's at column 4, which would require moving your 'I' to (4,5)? This way, we can keep the 'I's in a consistent row and column. What are your thoughts on this proposal?} \\
            }
        }
    }
    & & \\ \\

    \theutterance \stepcounter{utterance}  
    & & & \multicolumn{2}{p{0.3\linewidth}}{
        \cellcolor[rgb]{0.9,0.9,0.9}{
            \makecell[{{p{\linewidth}}}]{
                \texttt{\tiny{[GM$|$GM]}}
                \texttt{SAY: Placing the 'I's at row 5 seems like a reasonable approach, considering the current positions of the 'H's and 'I's. Since my 'I' is at (4,2) and your 'I' is at (1,5), we would need to move my 'I' to row 5. To maintain consistency, let's aim to place the 'I's at the same column as well. How about we place the 'I's at column 4, which would require moving your 'I' to (4,5)? This way, we can keep the 'I's in a consistent row and column. What are your thoughts on this proposal?} \\
            }
        }
    }
    & & \\ \\

    \theutterance \stepcounter{utterance}  
    & & \multicolumn{4}{p{0.6\linewidth}}{
        \cellcolor[rgb]{0.9,0.9,0.9}{
            \makecell[{{p{\linewidth}}}]{
                \texttt{\tiny{[GM$\rangle$P2]}}
                \texttt{Your message has been relayed to the other player.} \\
\\ 
\\ 
\texttt{You are currently playing round 12 of maximum 20.} \\
\texttt{You have collectively accumulated 9 of 12 penalties.} \\
\texttt{The other player sent this message:} \\
\\ 
\\ 
\texttt{"Placing the 'I's at row 5 seems like a reasonable approach, considering the current positions of the 'H's and 'I's. Since my 'I' is at (4,2) and your 'I' is at (1,5), we would need to move my 'I' to row 5. To maintain consistency, let's aim to place the 'I's at the same column as well. How about we place the 'I's at column 4, which would require moving your 'I' to (4,5)? This way, we can keep the 'I's in a consistent row and column. What are your thoughts on this proposal?"} \\
\\ 
\\ 
\texttt{What is your next command?} \\
            }
        }
    }
    & \\ \\

    \theutterance \stepcounter{utterance}  
    & & & \multicolumn{4}{p{0.6\linewidth}}{
        \cellcolor[rgb]{0.9,0.9,0.9}{
            \makecell[{{p{\linewidth}}}]{
                \texttt{\tiny{[GM$\langle$P2]}}
                \texttt{SAY: I agree with placing the 'I's at row 5 and column 4. Moving my 'I' to (4,5) is a good idea, as it would align with your 'I' at the same row and column. I'll make the move to (4,5). After that, we can focus on the 'W's and finalize their positions to complete the alignment.} \\
            }
        }
    }
     \\ \\

    \theutterance \stepcounter{utterance}  
    & & & \multicolumn{2}{p{0.3\linewidth}}{
        \cellcolor[rgb]{0.9,0.9,0.9}{
            \makecell[{{p{\linewidth}}}]{
                \texttt{\tiny{[GM$|$GM]}}
                \texttt{SAY: I agree with placing the 'I's at row 5 and column 4. Moving my 'I' to (4,5) is a good idea, as it would align with your 'I' at the same row and column. I'll make the move to (4,5). After that, we can focus on the 'W's and finalize their positions to complete the alignment.} \\
            }
        }
    }
    & & \\ \\

    \theutterance \stepcounter{utterance}  
    & & \multicolumn{4}{p{0.6\linewidth}}{
        \cellcolor[rgb]{0.9,0.9,0.9}{
            \makecell[{{p{\linewidth}}}]{
                \texttt{\tiny{[P1$\langle$GM]}}
                \texttt{Your message has been relayed to the other player.} \\
\\ 
\\ 
\texttt{You are currently playing round 12 of maximum 20.} \\
\texttt{You have collectively accumulated 9 of 12 penalties.} \\
\texttt{The other player sent this message:} \\
\\ 
\\ 
\texttt{"I agree with placing the 'I's at row 5 and column 4. Moving my 'I' to (4,5) is a good idea, as it would align with your 'I' at the same row and column. I'll make the move to (4,5). After that, we can focus on the 'W's and finalize their positions to complete the alignment."} \\
\\ 
\\ 
\texttt{What is your next command?} \\
            }
        }
    }
    & \\ \\

    \theutterance \stepcounter{utterance}  
    & \multicolumn{4}{p{0.6\linewidth}}{
        \cellcolor[rgb]{0.9,0.9,0.9}{
            \makecell[{{p{\linewidth}}}]{
                \texttt{\tiny{[P1$\rangle$GM]}}
                \texttt{SAY: I'm glad we're making progress. Before we move on to the 'W's, I'll wait for your confirmation that your 'I' has been moved to (4,5). Once that's done, we can discuss the 'W's. My 'I' needs to be moved to (4,5) as well, to match your 'I's new position. After we've aligned the 'I's, we can focus on finding a suitable position for the 'W's. Since we've been using a row{-}based approach for the 'H's and 'I's, perhaps we could consider a similar approach for the 'W's? What are your thoughts on placing the 'W's at a specific row or column that works for both our grids?} \\
            }
        }
    }
    & & \\ \\

    \theutterance \stepcounter{utterance}  
    & & & \multicolumn{2}{p{0.3\linewidth}}{
        \cellcolor[rgb]{0.9,0.9,0.9}{
            \makecell[{{p{\linewidth}}}]{
                \texttt{\tiny{[GM$|$GM]}}
                \texttt{SAY: I'm glad we're making progress. Before we move on to the 'W's, I'll wait for your confirmation that your 'I' has been moved to (4,5). Once that's done, we can discuss the 'W's. My 'I' needs to be moved to (4,5) as well, to match your 'I's new position. After we've aligned the 'I's, we can focus on finding a suitable position for the 'W's. Since we've been using a row{-}based approach for the 'H's and 'I's, perhaps we could consider a similar approach for the 'W's? What are your thoughts on placing the 'W's at a specific row or column that works for both our grids?} \\
            }
        }
    }
    & & \\ \\

    \theutterance \stepcounter{utterance}  
    & & \multicolumn{4}{p{0.6\linewidth}}{
        \cellcolor[rgb]{0.9,0.9,0.9}{
            \makecell[{{p{\linewidth}}}]{
                \texttt{\tiny{[GM$\rangle$P2]}}
                \texttt{Your message has been relayed to the other player.} \\
\\ 
\\ 
\texttt{You are currently playing round 13 of maximum 20.} \\
\texttt{You have collectively accumulated 9 of 12 penalties.} \\
\texttt{The other player sent this message:} \\
\\ 
\\ 
\texttt{"I'm glad we're making progress. Before we move on to the 'W's, I'll wait for your confirmation that your 'I' has been moved to (4,5). Once that's done, we can discuss the 'W's. My 'I' needs to be moved to (4,5) as well, to match your 'I's new position. After we've aligned the 'I's, we can focus on finding a suitable position for the 'W's. Since we've been using a row{-}based approach for the 'H's and 'I's, perhaps we could consider a similar approach for the 'W's? What are your thoughts on placing the 'W's at a specific row or column that works for both our grids?"} \\
\\ 
\\ 
\texttt{What is your next command?} \\
            }
        }
    }
    & \\ \\

    \theutterance \stepcounter{utterance}  
    & & & \multicolumn{4}{p{0.6\linewidth}}{
        \cellcolor[rgb]{0.9,0.9,0.9}{
            \makecell[{{p{\linewidth}}}]{
                \texttt{\tiny{[GM$\langle$P2]}}
                \texttt{MOVE: I, (4,5)} \\
            }
        }
    }
     \\ \\

    \theutterance \stepcounter{utterance}  
    & & & \multicolumn{2}{p{0.3\linewidth}}{
        \cellcolor[rgb]{0.9,0.9,0.9}{
            \makecell[{{p{\linewidth}}}]{
                \texttt{\tiny{[GM$|$GM]}}
                \texttt{MOVE: I, (4,5)} \\
            }
        }
    }
    & & \\ \\

    \theutterance \stepcounter{utterance}  
    & & & \multicolumn{2}{p{0.3\linewidth}}{
        \cellcolor[rgb]{0.9,0.9,0.9}{
            \makecell[{{p{\linewidth}}}]{
                \texttt{\tiny{[GM$|$GM]}}
                \texttt{Penalty: (4,5) is not empty, but contains '┤'.} \\
\\ 
\texttt{You have collectively accumulated 10 of 12 penalties.} \\
\texttt{Please try again!} \\
            }
        }
    }
    & & \\ \\

    \theutterance \stepcounter{utterance}  
    & & \multicolumn{4}{p{0.6\linewidth}}{
        \cellcolor[rgb]{0.9,0.9,0.9}{
            \makecell[{{p{\linewidth}}}]{
                \texttt{\tiny{[GM$\rangle$P2]}}
                \texttt{Penalty: (4,5) is not empty, but contains '┤'.} \\
\\ 
\texttt{You have collectively accumulated 10 of 12 penalties.} \\
\texttt{Please try again!} \\
            }
        }
    }
    & \\ \\

    \theutterance \stepcounter{utterance}  
    & & & \multicolumn{4}{p{0.6\linewidth}}{
        \cellcolor[rgb]{0.9,0.9,0.9}{
            \makecell[{{p{\linewidth}}}]{
                \texttt{\tiny{[GM$\langle$P2]}}
                \texttt{MOVE: I, (5,5)} \\
            }
        }
    }
     \\ \\

    \theutterance \stepcounter{utterance}  
    & & & \multicolumn{2}{p{0.3\linewidth}}{
        \cellcolor[rgb]{0.9,0.9,0.9}{
            \makecell[{{p{\linewidth}}}]{
                \texttt{\tiny{[GM$|$GM]}}
                \texttt{MOVE: I, (5,5)} \\
            }
        }
    }
    & & \\ \\

    \theutterance \stepcounter{utterance}  
    & & & \multicolumn{2}{p{0.3\linewidth}}{
        \cellcolor[rgb]{0.9,0.9,0.9}{
            \makecell[{{p{\linewidth}}}]{
                \texttt{\tiny{[GM$|$GM]}}
                \texttt{Penalty: (5,5) is not empty, but contains 'I'.} \\
\\ 
\texttt{You have collectively accumulated 11 of 12 penalties.} \\
\texttt{Please try again!} \\
            }
        }
    }
    & & \\ \\

    \theutterance \stepcounter{utterance}  
    & & \multicolumn{4}{p{0.6\linewidth}}{
        \cellcolor[rgb]{0.9,0.9,0.9}{
            \makecell[{{p{\linewidth}}}]{
                \texttt{\tiny{[GM$\rangle$P2]}}
                \texttt{Penalty: (5,5) is not empty, but contains 'I'.} \\
\\ 
\texttt{You have collectively accumulated 11 of 12 penalties.} \\
\texttt{Please try again!} \\
            }
        }
    }
    & \\ \\

    \theutterance \stepcounter{utterance}  
    & & & \multicolumn{4}{p{0.6\linewidth}}{
        \cellcolor[rgb]{0.9,0.9,0.9}{
            \makecell[{{p{\linewidth}}}]{
                \texttt{\tiny{[GM$\langle$P2]}}
                \texttt{MOVE: I, (3,5)} \\
            }
        }
    }
     \\ \\

    \theutterance \stepcounter{utterance}  
    & & & \multicolumn{2}{p{0.3\linewidth}}{
        \cellcolor[rgb]{0.9,0.9,0.9}{
            \makecell[{{p{\linewidth}}}]{
                \texttt{\tiny{[GM$|$GM]}}
                \texttt{MOVE: I, (3,5)} \\
            }
        }
    }
    & & \\ \\

    \theutterance \stepcounter{utterance}  
    & & & \multicolumn{2}{p{0.3\linewidth}}{
        \cellcolor[rgb]{0.9,0.9,0.9}{
            \makecell[{{p{\linewidth}}}]{
                \texttt{\tiny{[GM$|$GM]}}
                \texttt{Penalty: (3,5) is not empty, but contains '─'.} \\
\\ 
\texttt{You have collectively accumulated 12 of 12 penalties.} \\
\texttt{Please try again!} \\
            }
        }
    }
    & & \\ \\

    \theutterance \stepcounter{utterance}  
    & & \multicolumn{4}{p{0.6\linewidth}}{
        \cellcolor[rgb]{0.9,0.9,0.9}{
            \makecell[{{p{\linewidth}}}]{
                \texttt{\tiny{[GM$\rangle$P2]}}
                \texttt{Penalty: (3,5) is not empty, but contains '─'.} \\
\\ 
\texttt{You have collectively accumulated 12 of 12 penalties.} \\
\texttt{Please try again!} \\
            }
        }
    }
    & \\ \\

    \theutterance \stepcounter{utterance}  
    & & & \multicolumn{4}{p{0.6\linewidth}}{
        \cellcolor[rgb]{0.9,0.9,0.9}{
            \makecell[{{p{\linewidth}}}]{
                \texttt{\tiny{[GM$\langle$P2]}}
                \texttt{SAY: It seems I've reached the penalty limit. Before making another move, I'd like to discuss the 'I's position with you to avoid any further issues. Can we re{-}evaluate the grid and find a suitable position for the 'I's that works for both our grids?} \\
            }
        }
    }
     \\ \\

    \theutterance \stepcounter{utterance}  
    & & & \multicolumn{2}{p{0.3\linewidth}}{
        \cellcolor[rgb]{0.9,0.9,0.9}{
            \makecell[{{p{\linewidth}}}]{
                \texttt{\tiny{[GM$|$GM]}}
                \texttt{SAY: It seems I've reached the penalty limit. Before making another move, I'd like to discuss the 'I's position with you to avoid any further issues. Can we re{-}evaluate the grid and find a suitable position for the 'I's that works for both our grids?} \\
            }
        }
    }
    & & \\ \\

    \theutterance \stepcounter{utterance}  
    & & \multicolumn{4}{p{0.6\linewidth}}{
        \cellcolor[rgb]{0.9,0.9,0.9}{
            \makecell[{{p{\linewidth}}}]{
                \texttt{\tiny{[P1$\langle$GM]}}
                \texttt{Your message has been relayed to the other player.} \\
\\ 
\\ 
\texttt{You are currently playing round 13 of maximum 20.} \\
\texttt{You have collectively accumulated 12 of 12 penalties.} \\
\texttt{The other player sent this message:} \\
\\ 
\\ 
\texttt{"It seems I've reached the penalty limit. Before making another move, I'd like to discuss the 'I's position with you to avoid any further issues. Can we re{-}evaluate the grid and find a suitable position for the 'I's that works for both our grids?"} \\
\\ 
\\ 
\texttt{What is your next command?} \\
            }
        }
    }
    & \\ \\

    \theutterance \stepcounter{utterance}  
    & \multicolumn{4}{p{0.6\linewidth}}{
        \cellcolor[rgb]{0.9,0.9,0.9}{
            \makecell[{{p{\linewidth}}}]{
                \texttt{\tiny{[P1$\rangle$GM]}}
                \texttt{SAY: Given the penalty limit, it's crucial we make careful decisions. Let's re{-}evaluate the grid and focus on finding a suitable position for the 'I's. Since my 'I' is at (4,2) and your 'I' is at (1,5), and we previously discussed placing the 'I's at row 5 and column 4, I suggest we stick to that plan. To avoid further penalties, I'll move my 'I' to (4,5) to match the proposed position. Please confirm if this move aligns with your strategy, and then we can proceed with the 'W's. If we're both satisfied with the 'I's position, we can discuss the 'W's and try to finalize their positions to complete the alignment.} \\
            }
        }
    }
    & & \\ \\

    \theutterance \stepcounter{utterance}  
    & & & \multicolumn{2}{p{0.3\linewidth}}{
        \cellcolor[rgb]{0.9,0.9,0.9}{
            \makecell[{{p{\linewidth}}}]{
                \texttt{\tiny{[GM$|$GM]}}
                \texttt{SAY: Given the penalty limit, it's crucial we make careful decisions. Let's re{-}evaluate the grid and focus on finding a suitable position for the 'I's. Since my 'I' is at (4,2) and your 'I' is at (1,5), and we previously discussed placing the 'I's at row 5 and column 4, I suggest we stick to that plan. To avoid further penalties, I'll move my 'I' to (4,5) to match the proposed position. Please confirm if this move aligns with your strategy, and then we can proceed with the 'W's. If we're both satisfied with the 'I's position, we can discuss the 'W's and try to finalize their positions to complete the alignment.} \\
            }
        }
    }
    & & \\ \\

    \theutterance \stepcounter{utterance}  
    & & \multicolumn{4}{p{0.6\linewidth}}{
        \cellcolor[rgb]{0.9,0.9,0.9}{
            \makecell[{{p{\linewidth}}}]{
                \texttt{\tiny{[GM$\rangle$P2]}}
                \texttt{Your message has been relayed to the other player.} \\
\\ 
\\ 
\texttt{You are currently playing round 14 of maximum 20.} \\
\texttt{You have collectively accumulated 12 of 12 penalties.} \\
\texttt{The other player sent this message:} \\
\\ 
\\ 
\texttt{"Given the penalty limit, it's crucial we make careful decisions. Let's re{-}evaluate the grid and focus on finding a suitable position for the 'I's. Since my 'I' is at (4,2) and your 'I' is at (1,5), and we previously discussed placing the 'I's at row 5 and column 4, I suggest we stick to that plan. To avoid further penalties, I'll move my 'I' to (4,5) to match the proposed position. Please confirm if this move aligns with your strategy, and then we can proceed with the 'W's. If we're both satisfied with the 'I's position, we can discuss the 'W's and try to finalize their positions to complete the alignment."} \\
\\ 
\\ 
\texttt{What is your next command?} \\
            }
        }
    }
    & \\ \\

    \theutterance \stepcounter{utterance}  
    & & & \multicolumn{4}{p{0.6\linewidth}}{
        \cellcolor[rgb]{0.9,0.9,0.9}{
            \makecell[{{p{\linewidth}}}]{
                \texttt{\tiny{[GM$\langle$P2]}}
                \texttt{SAY: I confirm that moving your 'I' to (4,5) aligns with our previous discussion. Now that your 'I' is at (4,5), I will move my 'I' to a position that matches, but I need to find an empty space. Since we want the 'I's to be at the same row and column, I will try to move my 'I' to a position that is available. Let me check the grid again.} \\
            }
        }
    }
     \\ \\

    \theutterance \stepcounter{utterance}  
    & & & \multicolumn{2}{p{0.3\linewidth}}{
        \cellcolor[rgb]{0.9,0.9,0.9}{
            \makecell[{{p{\linewidth}}}]{
                \texttt{\tiny{[GM$|$GM]}}
                \texttt{SAY: I confirm that moving your 'I' to (4,5) aligns with our previous discussion. Now that your 'I' is at (4,5), I will move my 'I' to a position that matches, but I need to find an empty space. Since we want the 'I's to be at the same row and column, I will try to move my 'I' to a position that is available. Let me check the grid again.} \\
            }
        }
    }
    & & \\ \\

    \theutterance \stepcounter{utterance}  
    & & \multicolumn{4}{p{0.6\linewidth}}{
        \cellcolor[rgb]{0.9,0.9,0.9}{
            \makecell[{{p{\linewidth}}}]{
                \texttt{\tiny{[P1$\langle$GM]}}
                \texttt{Your message has been relayed to the other player.} \\
\\ 
\\ 
\texttt{You are currently playing round 14 of maximum 20.} \\
\texttt{You have collectively accumulated 12 of 12 penalties.} \\
\texttt{The other player sent this message:} \\
\\ 
\\ 
\texttt{"I confirm that moving your 'I' to (4,5) aligns with our previous discussion. Now that your 'I' is at (4,5), I will move my 'I' to a position that matches, but I need to find an empty space. Since we want the 'I's to be at the same row and column, I will try to move my 'I' to a position that is available. Let me check the grid again."} \\
\\ 
\\ 
\texttt{What is your next command?} \\
            }
        }
    }
    & \\ \\

    \theutterance \stepcounter{utterance}  
    & \multicolumn{4}{p{0.6\linewidth}}{
        \cellcolor[rgb]{0.9,0.9,0.9}{
            \makecell[{{p{\linewidth}}}]{
                \texttt{\tiny{[P1$\rangle$GM]}}
                \texttt{MOVE: I, (4, 5)} \\
            }
        }
    }
    & & \\ \\

    \theutterance \stepcounter{utterance}  
    & & & \multicolumn{2}{p{0.3\linewidth}}{
        \cellcolor[rgb]{0.9,0.9,0.9}{
            \makecell[{{p{\linewidth}}}]{
                \texttt{\tiny{[GM$|$GM]}}
                \texttt{MOVE: I, (4, 5)} \\
            }
        }
    }
    & & \\ \\

    \theutterance \stepcounter{utterance}  
    & & & \multicolumn{2}{p{0.3\linewidth}}{
        \cellcolor[rgb]{0.9,0.9,0.9}{
            \makecell[{{p{\linewidth}}}]{
                \texttt{\tiny{[GM$|$GM]}}
                \texttt{Penalty: (4,5) is not empty, but contains '┤'.} \\
\\ 
\texttt{You have collectively accumulated 13 of 12 penalties.} \\
\texttt{Please try again!} \\
            }
        }
    }
    & & \\ \\

    \theutterance \stepcounter{utterance}  
    & & & \multicolumn{2}{p{0.3\linewidth}}{
        \cellcolor[rgb]{0.9,0.9,0.9}{
            \makecell[{{p{\linewidth}}}]{
                \texttt{\tiny{[GM$|$GM]}}
                \texttt{Maximum number of penalties exceeded} \\
            }
        }
    }
    & & \\ \\

    \theutterance \stepcounter{utterance}  
    & & & \multicolumn{2}{p{0.3\linewidth}}{
        \cellcolor[rgb]{0.9,0.9,0.9}{
            \makecell[{{p{\linewidth}}}]{
                \texttt{\tiny{[GM$|$GM]}}
                \texttt{Initial states:} \\
\texttt{```} \\
\texttt{Player 1:} \\
\texttt{ 1234567} \\
\texttt{╔═╤═╤═╤═╗} \\
\texttt{║T│◌│◌│◌║ 1} \\
\texttt{╟─┼─┤I│◌║ 2} \\
\texttt{║W│◌│◌│H║ 3} \\
\texttt{║◌│◌├─┴─╢ 4} \\
\texttt{╟─┴─┤◌◌◌║ 5} \\
\texttt{║◌C◌└─┬─╢ 6} \\
\texttt{║◌◌◌◌◌│◌║ 7} \\
\texttt{╚═════╧═╝} \\
\\ 
\\ 
\texttt{Player 2:} \\
\texttt{ 1234567} \\
\texttt{╔═╤═╤═╤═╗} \\
\texttt{║◌│◌│◌│◌║ 1} \\
\texttt{╟─┼─┤◌│◌║ 2} \\
\texttt{║◌│W│◌│◌║ 3} \\
\texttt{║◌│C├─┴─╢ 4} \\
\texttt{╟─┴─┤I◌◌║ 5} \\
\texttt{║◌H◌└─┬─╢ 6} \\
\texttt{║◌◌T◌◌│◌║ 7} \\
\texttt{╚═════╧═╝} \\
\\ 
\texttt{```} \\
            }
        }
    }
    & & \\ \\

    \theutterance \stepcounter{utterance}  
    & & & \multicolumn{2}{p{0.3\linewidth}}{
        \cellcolor[rgb]{0.9,0.9,0.9}{
            \makecell[{{p{\linewidth}}}]{
                \texttt{\tiny{[GM$|$GM]}}
                \texttt{End states:} \\
\texttt{```} \\
\texttt{Player 1:} \\
\texttt{ 1234567} \\
\texttt{╔═╤═╤═╤═╗} \\
\texttt{║◌│◌│◌│◌║ 1} \\
\texttt{╟─┼─┤I│◌║ 2} \\
\texttt{║W│◌│◌│◌║ 3} \\
\texttt{║◌│◌├─┴─╢ 4} \\
\texttt{╟─┴─┤◌◌◌║ 5} \\
\texttt{║HC◌└─┬─╢ 6} \\
\texttt{║◌◌T◌◌│◌║ 7} \\
\texttt{╚═════╧═╝} \\
\\ 
\\ 
\texttt{Player 2:} \\
\texttt{ 1234567} \\
\texttt{╔═╤═╤═╤═╗} \\
\texttt{║◌│◌│◌│◌║ 1} \\
\texttt{╟─┼─┤◌│◌║ 2} \\
\texttt{║◌│W│◌│◌║ 3} \\
\texttt{║◌│◌├─┴─╢ 4} \\
\texttt{╟─┴─┤IC◌║ 5} \\
\texttt{║◌H◌└─┬─╢ 6} \\
\texttt{║◌◌T◌◌│◌║ 7} \\
\texttt{╚═════╧═╝} \\
\\ 
\texttt{```} \\
            }
        }
    }
    & & \\ \\

    \theutterance \stepcounter{utterance}  
    & & & \multicolumn{2}{p{0.3\linewidth}}{
        \cellcolor[rgb]{0.9,0.9,0.9}{
            \makecell[{{p{\linewidth}}}]{
                \texttt{\tiny{[GM$|$GM]}}
                \texttt{* success: False} \\
\texttt{* lose: True} \\
\texttt{* aborted: True} \\
\texttt{{-}{-}{-}{-}{-}{-}{-}} \\
\texttt{* Shifts: 2.00} \\
\texttt{* Max Shifts: 8.00} \\
\texttt{* Min Shifts: 4.00} \\
\texttt{* End Distance Sum: 10.12} \\
\texttt{* Init Distance Sum: 19.39} \\
\texttt{* Expected Distance Sum: 20.95} \\
\texttt{* Penalties: 13.00} \\
\texttt{* Max Penalties: 12.00} \\
\texttt{* Rounds: 14.00} \\
\texttt{* Max Rounds: 20.00} \\
\texttt{* Object Count: 5.00} \\
            }
        }
    }
    & & \\ \\

\end{supertabular}
}

\end{document}
